\documentclass[a4paper,10pt]{jsarticle}
\usepackage{amsmath,amssymb,amsthm}
\usepackage[dvipdfmx]{graphicx}
\usepackage{booktabs}
\usepackage{array}
\usepackage{url}

\usepackage[margin=25mm]{geometry}

\theoremstyle{definition}
\newtheorem{definition}{定義}[section]
\newtheorem{proposition}{命題}[section]
\newtheorem{lemma}{補題}[section]
\newtheorem{theorem}{定理}[section]
\newtheorem{conjecture}{予想}[section]
\newtheorem{remark}{注意}[section]
\makeatletter
\renewenvironment{proof}[1][\proofname]{\par
  \pushQED{\qed}%
  \normalfont \topsep6\p@\@plus6\p@\relax
  \trivlist
  \item[\hskip\labelsep\textbf{#1}]\ignorespaces
}{%
  \popQED\endtrivlist\@endpefalse
}
\renewcommand{\proofname}{証明}
\makeatother

\newcommand{\CVaR}{\mathrm{CVaR}}
\newcommand{\Wone}{W_1}
\expandafter\let\csname Finv\endcsname\undefined
\newcommand{\Finv}{F^{-1}}
\newcommand{\Fhinv}{\hat{F}^{-1}}
\newcommand{\R}{\mathbb{R}}
\newcommand{\E}{\mathbb{E}}
\newcommand{\Prob}{\mathbb{P}}
\newcommand{\dist}{\mathcal{D}_g}
\newcommand{\todo}[1]{\{#1\}}

\title{\vspace{-1cm}}
\author{}
\date{}

\begin{document}

%==================== 表紙 ====================
\begin{titlepage}
\begin{center}
\vspace*{3cm}
{\LARGE 卒業研究　中間レポート}\\[1.5cm]
{\Large 歪み関数重み付き1-Wasserstein距離に基づく\\[0.3cm]
quantile fraction学習と\\[0.3cm]
CVaR推定誤差の低減}\\[0.6cm]
{\large ---分布型強化学習FQFの歪みリスク尺度への適応---}\\[3cm]

\begin{tabular}{ll}
学籍番号 & 【学籍番号】\\[0.3cm]
氏名 & 【氏名】\\[0.3cm]
所属 & 【研究室名】\\[0.3cm]
指導教員 & 【指導教員名】\\[0.3cm]
提出日 & 2026年9月14日\\
\end{tabular}
\vfill
\end{center}
\end{titlepage}

\tableofcontents
\newpage

%==================== 概要 ====================
\section*{概要}
\addcontentsline{toc}{section}{概要}
分布型強化学習（distributional reinforcement learning）は，収益（リターン）の期待値ではなく分布そのものを学習することで，
リスク感応的な意思決定の基盤を与える枠組みである．
なかでも FQF（Fully Parameterized Quantile Function）\cite{yang2019fqf} は，分位点関数を評価するquantile fraction（累積確率軸 $[0, 1]$ の区切り位置）
$\tau_1<\dots<\tau_{N-1}$ を Fraction Proposal Network（FPN）により状態依存的に学習し，一様重みの1-Wasserstein距離 $\Wone(Z,\hat{Z})$（$Z$ は収益分布，$\hat{Z}$ はその近似分布）を最小化する．
しかし，金融におけるオプションのデルタヘッジのように，損益分布の下側裾を条件付きバリューアットリスク（CVaR）で評価したい応用では，
分布全体を一様に近似することが，CVaR推定という目的に対して効率的とは限らない．
本研究では，任意の歪みリスク尺度に対し限られた離散原子をどこに配置すればリスク尺度推定に有利かという観点から，
FPNの目的関数を歪みリスク尺度 $\rho_g$ の歪み関数 $g$ で重み付けした距離
$\Wone(\dist Z,\dist\hat{Z})=\int_0^1|\Finv(\tau)-\Fhinv(\tau)|g'(\tau)\,d\tau$ に置き換えることを提案する．
この置換により，(i) 各区間の原子を $g$-中央値にとることで母集団目的に関する境界勾配が FQF のquantile fractionに関する勾配の閉形式に係数 $g'(\tau_i)$ を乗じた形で与えられ，
分位点値ネットワークに plug-in した擬似勾配更新（stop-gradient 型）として実装できること，
(ii) 最適quantile fractionは歪み分布 $\dist Z$ の1次元最適量子化問題の解として特徴づけられること，
(iii) 既知の Lipschitz 評価の重み付き形として，歪みリスク尺度の推定誤差が提案損失によって上から抑えられることを示す．
数値実験Aでは，裾形状・非対称性の異なる3分布に対し，$\alpha N\ge2$ 領域における最小改善比 $R=42.3$ を達成し，$N=16$，$ \alpha=0.05$ では一様配置のCVaR推定誤差を改善比 $R=39.6$〜$97.9$ で低減した．
log--log 傾きの分析では，$\alpha=0.05$，$N=8$〜$1024$ の併合データにおける経験的傾きは，歪み最適配置で $-1.79$〜$-1.84$（おおむね二次減衰に整合），一様配置で $-0.72$〜$-1.27$（おおむね一次減衰）であり，この差の機序は下側裾の特異性に対する適応配置の回避にある（注意\ref{rem:midpoint}）．改善比は $N$ とともに増大し（$N=1024$ で $R=1.1\times10^3$〜$4.2\times10^3$），一般的な漸近次数の証明は今後の課題である．
点数換算値（N-equivalence）の観察では，一様配置が提案法と同等の精度を得るにはガウス混合で約695点を要し，残る2分布では探索上限の1024点でも達しなかった．
自作 toy 環境 RiskyChain-v0 での実験Bでは，$g$-一様配置で初期化した提案法（B1g）のCVaR推定MAE（平均絶対誤差，mean absolute error）が
FQF（B0）の約7分の1（$7.02\to0.96$）に減少し，下側裾への配置（平均14.06個）を確認した一方，
リスク感応IQN（B2）がさらに小さい誤差（$0.36$）を示し，提案法の一般的優位性は未確認である．
さらに，quantile fraction を固定した ablation では，$g$-一様配置と $g$-中央値原子を用いる B3g が $0.74\pm0.33$ の MAE を示した．3シードのため B1g との優劣は確定できないが，toy 環境では改善の主因が FPN の状態依存最適化よりも，歪み重みに整合した理論配置の採用にあることが示唆された（FPN は状態依存化の一実装であり，その追加的寄与は状態ごとに最適配置が異なる環境で検証する）．
なおこの比較の参照値は各手法の greedy 方策下の分布に対するものであるため，共通方策下の交差評価も追加した．
その結果，all-RISKY 方策下の重い裾の参照では全評価器の誤差が $4.3$〜$7.0$ に達し，推定対象の難易度が誤差を支配することが示唆された．ただし本評価は学習済みネットワーク出力と固定方策参照値の距離（学習時の継続方策の不一致を含む）を測るものであり，交絡を完全には分離できない．厳密な分離には方策評価モードでの再学習が必要であるため，これを今後の課題とする（第\ref{subsec:cross}節）．
本レポートでは，これまでに理解した理論的基礎と導出，実験結果，および今後デルタヘッジ環境へ展開する計画を報告する．
%==================== 1 はじめに ====================
\section{はじめに}
\label{sec:intro}

\subsection{研究動機}
強化学習における価値関数 $Q^\pi(s,a)=\E[Z^\pi(s,a)]$ は，収益 $Z^\pi(s,a)=\sum_t\gamma^t R_t$ の期待値のみを扱う．
しかし，金融リスク管理をはじめ多くの意思決定問題では，収益の「平均」よりも「稀に発生する大きな損失」を正確に把握することが重要である．
たとえばオプションを売却した投資家は，原資産を売買してデルタヘッジを行うが，離散的なリヘッジと取引コストのもとでは
ヘッジ損益の分布は非対称で，厚い下側裾をもつ裾の重い分布（ヘビーテイル分布：裾が正規分布より重く，稀だが巨大な損失が無視できない確率で発生する型の分布）となる．
このとき投資家が重視すべきは分散ではなく，最悪側の累積確率 $\alpha$ における平均損失である．本報告では以下，損益 $Z$ を基準としてこの量を「下側 CVaR$_\alpha$ の\textbf{最大化}」という規約で扱う（第\ref{sec:drm}節の符号規約を参照）\cite{rockafellar2000cvar}．

分布型強化学習 \cite{bellemare2017c51,dabney2018qrdqn,dabney2018iqn,yang2019fqf} は収益分布 $Z^\pi$ 全体を学習するため，
学習した分布から $\CVaR_\alpha$ を計算してリスク回避的な方策を導くことが原理的に可能である．
実際，IQN \cite{dabney2018iqn} では行動選択時の分位点サンプリングを歪めることで CVaR などのリスク尺度に基づく行動選択を行い，
Moghimi and Ku \cite{moghimi2025beyond} は静的スペクトルリスク尺度を最適化する分布型RLを提案している．
しかし，これらの手法において「分布をどの点で近似するか」は，リスク尺度とは無関係に，
分布全体を一様に近似するよう定められている．
$\CVaR_{0.05}$ を評価するために必要なのは分位点関数 $\Finv(\tau)$ の $\tau\in[0,0.05]$ における形状であるが，
$N=16$ 個の一様quantile fractionのうち，内部の $\tau_i=i/16$（$i=1,\dots,N-1$）でこの区間に入るものは1つもない．
これが本研究の出発点である．\textbf{収益分布の下側裾（損失側の裾）を正確に評価したいなら，
近似の「解像度」を下側裾に集中させるべきである}．

\subsection{本研究の位置づけと貢献}
FQF \cite{yang2019fqf} はquantile fraction（累積確率軸 $[0, 1]$ の区切り位置）$\tau_i$ 自体を状態依存的に学習する代表的な分布型強化学習手法であり，
quantile fractionを一様重み $\Wone$ の最小化により決めている．
本研究はこの FQF の Fraction Proposal Network の目的関数を，歪み関数 $g$ で重み付けした $\Wone(\dist Z,\dist\hat{Z})$ に置き換える．
本研究は，CVaR最適方策を学習する新しい強化学習アルゴリズムの提案というよりは，
「任意の歪みリスク尺度に対し，限られた離散原子をどこに配置すればリスク尺度推定に有利か」という表現設計・目的指向量子化の研究として位置づけられる．
現時点での貢献は以下の通りである．
\begin{enumerate}
  \item \textbf{理論・アルゴリズム}：各区間の最適原子が $g$-中央値で与えられること（命題\ref{prop:gmedian}），および母集団目的に対する境界勾配が FQF の閉形式に係数 $g'(\tau_i)$ を乗じた形になること（命題\ref{prop:grad}）を導出し，これを分位点値ネットワークに plug-in した擬似勾配更新（stop-gradient 型）として実装可能な閉形式を与えた（第\ref{sec:proposed}節）．あわせて，既知の Lipschitz 評価の重み付き形として $|\rho_g(Z)-\rho_g(\hat Z)|\le \Wone(\dist Z,\dist \hat Z)$ が成り立つことを示した（第\ref{sec:drm}節）．
  \item \textbf{最適量子化としての解釈}：FPN の目的関数の最適化が，歪み分布 $\dist Z$ の1次元 $L^1$ 最適量子化問題 \cite{graf2000quantization} と数学的に同値であることを明らかにし，quantile fraction の再配置を目的指向量子化として位置づけた（第\ref{sec:proposed}節）．
  \item \textbf{実験A}：真の分位点関数が解析的に与えられた1次元設定で，$\alpha N\ge2$ 領域における最小改善比 $R=42.3$（$N=16, \alpha=0.05$ では $R=39.6$〜$97.9$）を達成し，原理的有効性を確認した（第\ref{sec:exp}節）．
  \item \textbf{実験B}：ヘビーテイル損失を持つ toy MDP 上で実装し，$g$-一様初期化を施した提案法（B1g）において CVaR 推定 MAE が FQF（B0）の約7分の1（$7.02\to0.96$）に低減することを確認した（第\ref{sec:exp}節）．ただしリスク感応 IQN がさらに小さい誤差（$0.36$）を示し，提案法の一般的優位性は未確認である．さらに固定配置の ablation（B3g）により，toy 環境での改善の主因がリスク尺度に整合した配置と原子選択にあることを示し，FPN による状態依存化の追加的有効性は状態依存環境で検証する課題とした．また，greedy 方策の違いによる交絡を共通方策下の交差評価で検討した結果，all-RISKY 分布（重い裾）の参照では全評価器の誤差が大きく推定対象の難易度が支配的であることが示唆された一方，本評価は学習時の継続方策の不一致を含むため，表現精度のみの厳密な分離は方策評価モードでの再学習を要する（第\ref{subsec:cross}，第\ref{sec:future}節）．
\end{enumerate}
本レポートは中間報告であり，デルタヘッジ環境での実験と誤差評価の理論的精密化は今後の課題である（第\ref{sec:future}節）．
%==================== 2 分布型強化学習の基礎 ====================
\section{分布型強化学習の基礎}
\label{sec:drl}

\subsection{収益分布と分布ベルマン作用素}
マルコフ決定過程 $(\mathcal{S},\mathcal{A},P,R,\gamma)$ と方策 $\pi$ に対し，収益を確率変数
\begin{equation}
  Z^\pi(s,a)=\sum_{t=0}^{\infty}\gamma^t R_t,\qquad s_0=s,\ a_0=a,\ a_t\sim\pi(\cdot|s_t)
\end{equation}
と定める．$Q^\pi(s,a)=\E[Z^\pi(s,a)]$ である．Bellemare ら \cite{bellemare2017c51} は分布ベルマン作用素
\begin{equation}
  (\mathcal{T}^\pi Z)(s,a)\overset{D}{=}R(s,a)+\gamma Z(S',A'),\qquad S'\sim P(\cdot|s,a),\ A'\sim\pi(\cdot|S')
\end{equation}
を導入し（$\overset{D}{=}$ は分布として等しいことを表す），これが最大形式の $p$-Wasserstein 距離 $\bar d_p(Z_1,Z_2)=\sup_{s,a}W_p(Z_1(s,a),Z_2(s,a))$（$\sup$ は全状態・行動対に対する上限）に関して
$\gamma$-縮小写像であることを示した．
ここで1次元分布 $Z_1,Z_2$ の間の $p$-Wasserstein距離は分位点関数（累積分布関数の一般化逆関数）$F^{-1}_{Z}(\tau)=\inf\{z:F_Z(z)\ge\tau\}$ を用いて
\begin{equation}
  W_p(Z_1,Z_2)=\left(\int_0^1\left|F_{Z_1}^{-1}(\tau)-F_{Z_2}^{-1}(\tau)\right|^p d\tau\right)^{1/p}
  \label{eq:wp}
\end{equation}
と表される．本研究で中心となるのは $p=1$ の場合
\begin{equation}
  \Wone(Z,\hat Z)=\int_0^1\left|\Finv(\tau)-\Fhinv(\tau)\right|d\tau
  \label{eq:w1}
\end{equation}
であり，これを以後「一様重み $\Wone$」と呼ぶ．積分測度が $[0,1]$ 上の一様測度 $d\tau$ であることを強調するためである．

\subsection{C51：固定サポート上のカテゴリカル分布}
C51 \cite{bellemare2017c51} は収益のサポート（分布が確率質量をもつ値の範囲）$[V_{\min},V_{\max}]$ 上に，等間隔の $N=51$ 個の離散的な点
$z_1<\dots<z_N$ を固定し，各点の確率 $p_i(s,a)$ をソフトマックスで出力する．
このように確率質量を置く離散的な点を\textbf{原子（アトム，atom）}と呼ぶ．
ベルマン更新後の分布 $R+\gamma Z'$ の原子は元のサポート上に乗らないため，隣接原子に確率質量を線形分配する射影 $\Phi$ を施し，
射影後の分布との交差エントロピーを最小化する．
理論上は Wasserstein 距離での縮小性が示されているのに対し，実際の損失が KL であるという理論と実装の乖離があった．
またサポート範囲を事前に設定する必要があり，その範囲外の収益が生じる場合には端点へのクリッピング誤差が発生する．
このため，ヘビーテイルの収益分布ではサポート範囲の設定に注意を要する．

\subsection{QR-DQN：分位点回帰による固定quantile fraction表現}
QR-DQN \cite{dabney2018qrdqn} は原子の位置ではなく確率側を固定する．
すなわち，各原子に均等な確率 $1/N$ を割り当て，ネットワークがその位置 $\theta_i(s,a)$ を出力することで，$Z(s,a)\approx\frac1N\sum_{i=1}^N\delta_{\theta_i(s,a)}$ と近似する．ここで $\delta_\theta$ は点 $\theta$ に確率質量1を置く離散分布（ディラック測度）である．
Dabney ら \cite{dabney2018qrdqn} は，等確率 $1/N$ の $N$ 点分布のうち $\Wone(Z,\hat Z)$ を最小化する離散分布の原子が，分位点中点 $\hat\tau_i=(2i-1)/(2N)$ における分位点値
$\theta_i=\Finv(\hat\tau_i)$ で与えられることを示した．QR-DQN はこの原子位置を分位点回帰で学習する．
分位点 $\tau$ に対する分位点回帰損失は，指示関数 $\mathbb{I}$ を用いて
\begin{equation}
  \rho_\tau(u)=u\left(\tau-\mathbb{I}\{u<0\}\right),\qquad u=\text{（目標）}-\text{（予測）}
\end{equation}
であり，その期待値の最小点が $\tau$-分位点になる．
誤差ゼロ付近での非微分性（滑らかでないこと）に起因する学習の不安定性を避けるため，実装では Huber 損失 $\mathcal{L}_\kappa$ を組み合わせた
\begin{equation}
  \rho^\kappa_\tau(u)=\left|\tau-\mathbb{I}\{u<0\}\right|\frac{\mathcal{L}_\kappa(u)}{\kappa},\qquad
  \mathcal{L}_\kappa(u)=\begin{cases}\tfrac12u^2 & |u|\le\kappa\\ \kappa(|u|-\tfrac12\kappa)& \text{それ以外}\end{cases}
  \label{eq:huber-quantile}
\end{equation}
を用いる．目標分布の原子 $\theta'_j$ と予測原子 $\theta_i$ のすべての対について
\begin{equation}
  \mathcal{L}_{\mathrm{QR}}=\frac1N\sum_{i=1}^N\sum_{j=1}^N\rho^\kappa_{\hat\tau_i}\!\left(r+\gamma\theta'_j-\theta_i\right)
  \label{eq:qr-loss}
\end{equation}
を最小化する．この定式化により損失が Wasserstein 距離の最小化と整合的になり，サポートの事前固定も不要になった．
なお，式\eqref{eq:qr-loss}は QR-DQN や IQN における標準形（等重み $1/N$ による平均）を示しているが，第\ref{sec:proposed}節で述べる FQF 系ではターゲット原子が区間幅に応じて重み付けされる点に留意されたい．
本研究の観点で重要なのは，QR-DQN の近似能力が「quantile fractionが一様であること」に縛られている点である．

\subsection{IQN：暗黙的分位点ネットワーク}
IQN \cite{dabney2018iqn} はquantile fraction $\tau$ をネットワークの入力とし，
$\tau\sim U(0,1)$ をサンプリングし，$Z_\tau(s,a)=f(\psi(s),\phi(\tau))_a$ を出力する．
$\tau$ の埋め込みにはコサイン基底 $\phi_j(\tau)=\mathrm{ReLU}\!\left(\sum_{k=0}^{n-1}\cos(\pi k\tau)w_{jk}+b_j\right)$ を用い，
状態特徴 $\psi(s)$ と要素積で結合する．損失は式\eqref{eq:qr-loss}と同形で，$\tau_i,\tau'_j$ を独立にサンプリングして計算する．
IQN の特筆すべき性質は，行動選択時に $\tau$ のサンプリング分布を歪めることで，
\begin{equation}
  \pi_\beta(s)=\arg\max_a\ \E_{\tau\sim\beta}\left[Z_\tau(s,a)\right]
\end{equation}
としてリスク感応的な方策を実現できる点である．ここで $\beta$ は行動選択時の $\tau$ のサンプリング分布（リスクプロファイル）である．
とくに $\tau\sim U(0,\alpha)$ とすれば $\E_{\tau\sim U(0,\alpha)}[Z_\tau]$ は $\CVaR_\alpha$ の推定量となる．
これは第\ref{sec:drm}節で述べる歪みリスク尺度と本質的に同じ考え方であり，本研究の実験Bではベースライン B2 として採用する．
ただし IQN において「どの $\tau$ を用いて学習するか」と「どの $\tau$ を用いて行動選択するか」は別々に決められており，
IQN 原論文の設定としては学習時の $\tau$ の配置はリスク尺度から独立に一様である（ただし実装上は学習時のサンプリング分布を変更することも可能であり，第\ref{sec:exp}節で述べる対照 B2p はこの設定に対応する）．

\subsection{FQF：quantile fractionの完全パラメータ化}
\label{subsec:fqf}
FQF \cite{yang2019fqf} は QR-DQN の固定quantile fraction $\tau_i=i/N$ を，状態 $s$ の関数として学習する．
ネットワークは2つの部分からなる．
\begin{itemize}
  \item \textbf{Fraction Proposal Network（FPN）} $\omega_1$：状態特徴 $\psi(s)$ から $N$ 個の値を出力し，ソフトマックスで確率ベクトル $p\in\Delta^{N-1}$（$\Delta^{N-1}$ は各成分が非負で総和1のベクトルの集合）に変換したのち累積和をとって
  $\tau_0=0<\tau_1<\dots<\tau_{N-1}<\tau_N=1$ を得る．
  \item \textbf{Quantile Value Network} $\omega_2$：IQN と同様にquantile fraction $\hat\tau_i$ を入力として分位点値 $\Finv_{\omega_2}(\hat\tau_i)$ を出力する．
\end{itemize}
ここで $\hat\tau_i=(\tau_i+\tau_{i+1})/2$ は区間 $[\tau_i,\tau_{i+1}]$ の中点である．
近似分布は原子 $\Finv(\hat\tau_i)$ と重み $\tau_{i+1}-\tau_i$ をもつ離散分布
\begin{equation}
  \hat Z_N=\sum_{i=0}^{N-1}(\tau_{i+1}-\tau_i)\,\delta_{\Finv(\hat\tau_i)}
  \label{eq:zhat}
\end{equation}
であり，その分位点関数は階段関数 $\Fhinv(u)=\Finv(\hat\tau_i)$，$u\in(\tau_i,\tau_{i+1}]$ となる．
分位点値ネットワークは式\eqref{eq:qr-loss}の分位点回帰損失で，FPN は $\Wone(Z,\hat Z_N)$ の最小化でそれぞれ学習する．

\paragraph{補題1（一様重み $\Wone$ に対する最適原子）．}
Yang ら \cite{yang2019fqf} の Lemma 1 は，任意の連続分位点関数 $\Finv$ に対し，
quantile fraction $\tau_0<\dots<\tau_N$ を固定したとき，式\eqref{eq:w1}の一様重み $\Wone(Z,\hat Z_N)$ を最小化する原子の位置が，
各区間の中点 $\hat\tau_i=(\tau_i+\tau_{i+1})/2$ における分位点値 $\Finv(\hat\tau_i)$ で与えられることを示している．
これは次のように理解できる．区間 $i$ の寄与
\begin{equation}
  \ell_i(q)=\int_{\tau_i}^{\tau_{i+1}}\left|\Finv(u)-q\right|du
\end{equation}
を原子 $q$ で微分すると，$\Finv$ の単調性より $\partial\ell_i/\partial q=\Prob(u\in[\tau_i,\tau_{i+1}]:\Finv(u)<q)-\Prob(u\in[\tau_i,\tau_{i+1}]:\Finv(u)\ge q)$ であり，
これが零になるのは $q=\Finv(m)$，$m$ が区間上の一様測度の中央値，すなわち $m=\hat\tau_i$ のときである．
第\ref{sec:proposed}節ではこの「区間上の測度の中央値」という見方が，歪み重みへの拡張の鍵となる．

\paragraph{quantile fractionに関する勾配の閉形式．}
FPN は $\Wone(Z,\hat Z_N)$ の $\tau_i$（$i=1,\dots,N-1$）に関する勾配で学習される．
Yang ら \cite{yang2019fqf} は Proposition 1 で，これが
\begin{equation}
  \frac{\partial \Wone}{\partial\tau_i}=2\Finv(\tau_i)-\Finv(\hat\tau_i)-\Finv(\hat\tau_{i-1})
  \label{eq:fqf-grad}
\end{equation}
と閉形式で書けることを示した．導出を確認しておく．
$\Wone=\sum_{j=0}^{N-1}\ell_j(\Finv(\hat\tau_j))$ において $\tau_i$ に依存する項は区間 $i-1$ と区間 $i$ のみである．
$\tau_i$ は区間 $i-1$ の積分上限かつ区間 $i$ の積分下限として現れるため，Leibnizの積分則により
\begin{equation}
  \left|\Finv(\tau_i)-\Finv(\hat\tau_{i-1})\right|-\left|\Finv(\tau_i)-\Finv(\hat\tau_i)\right|
  =\left(\Finv(\tau_i)-\Finv(\hat\tau_{i-1})\right)-\left(\Finv(\hat\tau_i)-\Finv(\tau_i)\right)
\end{equation}
を得る．ここで $\hat\tau_{i-1}\le\tau_i\le\hat\tau_i$ と $\Finv$ の単調性を用いて絶対値を外した．
一方，原子 $\Finv(\hat\tau_{i-1}),\Finv(\hat\tau_i)$ を通じた依存（$\partial\hat\tau_i/\partial\tau_i=1/2$）は，
補題1により各原子の位置 $q_j$ が区間損失 $\ell_j$ の最小点であるため $\partial\ell_j/\partial q=0$ となり消える．
両者を合わせて式\eqref{eq:fqf-grad}が得られる．
直感的には，$\Finv(\tau_i)$ が隣接する2つの原子の中点にあるとき勾配が零になる．
すなわち一様重み $\Wone$ の停留点では「境界の分位点値＝隣接原子の中点」「原子＝区間中点の分位点値」という2条件が同時に成り立ち，
これは1次元 $L^1$ 最適量子化（Lloyd の条件；$k$-means 法の更新規則と同型の固定点条件）そのものである \cite{graf2000quantization}．

\subsection{系譜のまとめ}
表\ref{tab:lineage}に4手法を整理する．C51からFQFへの発展は「値側の固定」から「確率側の固定」，そして「確率側の学習」への移行である．
ただし FQF においてもquantile fractionの学習基準は一様重み $\Wone$ であり，リスク尺度は反映されていない．
本研究はこの最後の欄を埋めるものである．

\begin{table}[htbp]
\centering
\caption{分布型強化学習の系譜．最右列が本研究で置き換える部分．}
\label{tab:lineage}
\small
\setlength{\tabcolsep}{2.5pt}
\begin{tabular}{lllll}
\toprule
手法 & 表現 & 固定するもの & 損失 & quantile fractionの基準\\
\midrule
C51 \cite{bellemare2017c51} & カテゴリカル & 原子位置 & 射影後の交差エントロピー & ---\\
QR-DQN \cite{dabney2018qrdqn} & 分位点（等確率） & quantile fraction $i/N$ & 分位点回帰 & 一様（固定）\\
IQN \cite{dabney2018iqn} & 暗黙的分位点関数 & なし（サンプリング） & 分位点回帰 & 一様（サンプリング）\\
FQF \cite{yang2019fqf} & 分位点（可変確率） & なし & 分位点回帰＋$\Wone$ & 一様重み $\Wone$ 最小化\\
\textbf{提案} & 分位点（可変確率） & なし & 分位点回帰＋$\Wone(\dist Z,\dist\hat Z)$ & \textbf{歪み重み $\Wone$ 最小化}\\
\bottomrule
\end{tabular}
\end{table}

%==================== 3 歪みリスク尺度 ====================
\section{歪みリスク尺度}
\label{sec:drm}

\subsection{定義}
以下，$Z$ は損益（大きいほど良い）を表す実数値確率変数とし，リスクは下側裾にあるとする．
したがって本研究で「リスク尺度」と呼ぶ $\rho_g(Z)$ は，損失 $L=-Z$ に対して最小化されるリスク尺度というよりも，\textbf{収益 $Z$ に対する歪み重み付き評価（risk-adjusted value）}であり，下側 CVaR の\textbf{最大化}を目的とする規約で一貫させる（損失側・最小化の規約と混同しないこと）．
\begin{definition}[歪み関数と歪みリスク尺度]
非減少で $g(0)=0$，$g(1)=1$ を満たす関数 $g:[0,1]\to[0,1]$ を歪み関数と呼ぶ．
$g$ が絶対連続で導関数 $g'$ をもつとき，歪みリスク尺度を
\begin{equation}
  \rho_g(Z)=\int_0^1\Finv(\tau)\,dg(\tau)=\int_0^1\Finv(\tau)\,g'(\tau)\,d\tau
  \label{eq:drm}
\end{equation}
と定義する．
\end{definition}
なお，本研究で用いる歪み関数は，後述の $g_\alpha$，$g_{\alpha,\lambda}$，$g_\eta$ を含めいずれも区分的になめらかであり，$g'$ は区分定数または連続である．以下ではこの正則性を仮定すれば十分である．$g(\tau)=\tau$ のとき $\rho_g(Z)=\E[Z]$ である．
$g$ が凹関数のとき $g'$ は非増加であり，小さい $\tau$（悪い結果）に大きい重みが乗るためリスク回避的な評価となる．
この形の尺度は Yaari の双対理論 \cite{yaari1987dual} に由来し，期待効用理論が「値」を効用関数で変換するのに対し，
「確率」を歪み関数で変換する．
Moghimi and Ku \cite{moghimi2025beyond} が扱うスペクトルリスク尺度 $\int_0^1\Finv(\tau)\phi(\tau)d\tau$ は，
リスクスペクトル $\phi=g'$ と対応させれば式\eqref{eq:drm}と同じ族である．

\subsection{CVaR とその分位点表現}
\label{subsec:cvar}
\begin{definition}[CVaR]
水準 $\alpha\in(0,1)$ に対し，損益 $Z$ の条件付きバリューアットリスクを
\begin{equation}
  \CVaR_\alpha(Z)=\frac1\alpha\int_0^\alpha\Finv(\tau)\,d\tau
  \label{eq:cvar}
\end{equation}
と定義する．
\end{definition}
式\eqref{eq:cvar}の分位点積分表現（一般化逆関数 $\Finv$ による）を本報告の基準定義とする．連続分布では $\CVaR_\alpha(Z)=\E[Z\mid Z\le\Finv(\alpha)]$，すなわち最悪側の累積確率 $\alpha$ における損益の平均と一致するが，この条件付き期待値表現は連続分布に限った解釈であり，点質量を含む分布（実験Bの all-SAFE 方策下など）では式\eqref{eq:cvar}または Rockafellar--Uryasev 変分表現を基準に用いる．
Rockafellar and Uryasev \cite{rockafellar2000cvar} は
\begin{equation}
  \CVaR_\alpha(Z)=\max_{\eta\in\R}\left\{\eta-\frac1\alpha\E\left[(\eta-Z)^+\right]\right\}
\end{equation}
という変分表現を与え，CVaR最適化が凸計画に帰着することを示した（損失を正で扱う原論文と符号規約が逆である点に注意する）．ここで $\eta$ は閾値を表す補助変数であり，この最大問題の最適解は $\Finv(\alpha)$ となる．また $(\eta-Z)^+=\max(\eta-Z,0)$ は正位部分を表す．
式\eqref{eq:cvar}は歪み関数
\begin{equation}
  g_\alpha(\tau)=\min\left(\tau/\alpha,\,1\right),\qquad
  g'_\alpha(\tau)=\begin{cases}1/\alpha & 0\le\tau<\alpha\\ 0 & \alpha<\tau\le 1\end{cases}
  \label{eq:g-cvar}
\end{equation}
に対する歪みリスク尺度 $\rho_{g_\alpha}(Z)$ に一致する．
$g'_\alpha$ は $[0,\alpha)$ で一定値 $1/\alpha$，$(\alpha,1]$ で零であり，
「分位点関数の $[0,\alpha]$ 部分だけを，一様測度の $1/\alpha$ 倍の密度で見ている」ことが分かる．
この観察が，quantile fractionを $[0,\alpha]$ に集中させるべきという第\ref{sec:intro}節の直感の数学的表現である．

\subsection{歪み分布と推定誤差の上界}
\label{subsec:bound}
$\hat Z$ を $Z$ の近似分布（たとえば式\eqref{eq:zhat}の $\hat Z_N$）とする．
$\hat Z$ から計算した $\rho_g(\hat Z)$ が真の $\rho_g(Z)$ からどれだけ離れうるかを評価する．

\begin{definition}[歪み分布]
\label{def:distorted}
絶対連続な歪み関数 $g$ に対し，累積分布関数が $g$ である $[0,1]$ 上の確率変数 $U_g$ をとる．
$Z$ の $g$ による歪み分布 $\dist Z$（または $\mathcal{D}_g[Z]$）を，確率変数 $\Finv(U_g)$ の分布と定める（以下，$g^{-1}$ は $g$ の一般化逆関数を表す）．
なお記号 $\dist Z$ は，値 $Z$ に関数 $g$ を作用させる通常の押し出し測度ではなく，
分位レベルの分布を $g$ により歪めて得られる分布を表す本稿の略記である．
\end{definition}
$\Prob(\Finv(U_g)\le z)=\Prob(U_g\le F(z))=g(F(z))$ であるから，$\dist Z$ の累積分布関数は $g\circ F_Z$ であり，
適切な端点の取扱いの下で，その分位点関数はほとんど至る所
$F^{-1}_{\dist Z}(v)=\Finv(g^{-1}(v))$ と表せる．
$Z$ が確率密度関数 $f$ をもつとき，$\dist Z$ の確率密度関数は $g'(F(z))f(z)$ となり，元の確率密度を重み $g'(F(z))$ で再重み付けした分布に他ならない．
特に
\begin{equation}
  \E[\dist Z]=\int_0^1\Finv(g^{-1}(v))\,dv=\int_0^1\Finv(\tau)\,g'(\tau)\,d\tau=\rho_g(Z)
  \label{eq:drm-as-mean}
\end{equation}
であり（$v=g(\tau)$ と置換），\textbf{歪みリスク尺度は歪み分布の平均}である．

\begin{proposition}[歪み重み付き $\Wone$ 上界]
\label{prop:bound}
$g$ を絶対連続な歪み関数とし，$Z,\hat Z$ は $\Finv,\Fhinv$ が測度 $g'(\tau)\,d\tau$ に関して可積分である（すなわち右辺が有限となる．$g'$ が有界な場合，$\E|Z|<\infty$，$\E|\hat Z|<\infty$ が十分条件であり，本実験の Student-$t$（自由度3）もこれを満たす．なお $\hat Z_N$ は有限離散分布なので $\hat Z$ 側の可積分性は自明である）とする．このとき任意の $Z,\hat Z$ に対して
\begin{equation}
  \left|\rho_g(Z)-\rho_g(\hat Z)\right|
  \le\int_0^1\left|\Finv(\tau)-\Fhinv(\tau)\right|g'(\tau)\,d\tau
  =\Wone(\dist Z,\dist\hat Z).
  \label{eq:bound}
\end{equation}
\end{proposition}
\begin{proof}
第1の不等式は式\eqref{eq:drm}から直接
\begin{equation*}
  \left|\rho_g(Z)-\rho_g(\hat Z)\right|=\left|\int_0^1\left(\Finv(\tau)-\Fhinv(\tau)\right)g'(\tau)\,d\tau\right|
  \le\int_0^1\left|\Finv(\tau)-\Fhinv(\tau)\right|g'(\tau)\,d\tau
\end{equation*}
として得られる（$g'\ge0$ を用いた）．第2の等式は，式\eqref{eq:wp}を $\dist Z,\dist\hat Z$ に適用し
\begin{equation*}
  \Wone(\dist Z,\dist\hat Z)=\int_0^1\left|\Finv(g^{-1}(v))-\Fhinv(g^{-1}(v))\right|dv
\end{equation*}
において $v=g(\tau)$，$dv=g'(\tau)d\tau$ と置換すれば従う．
\end{proof}
本命題は，確率測度の差異に対するリスク尺度の Lipschitz 連続性に関する評価（Pichler \cite{pichler2013evaluations}）の重み付き形に相当し，初等的な三角不等式から直ちに従う自明な上界であるが，FQF の目的関数をリスク尺度に応じて設計するための直接的な動機を与える．

\begin{remark}
式\eqref{eq:bound}の右辺は一様重み $\Wone(Z,\hat Z)$ の被積分関数に重み $g'(\tau)$ を掛けたものである．
CVaR の場合 $g'_\alpha=\alpha^{-1}\mathbb{I}_{[0,\alpha)}$ であるから
\begin{equation}
  \left|\CVaR_\alpha(Z)-\CVaR_\alpha(\hat Z)\right|\le\frac1\alpha\int_0^\alpha\left|\Finv(\tau)-\Fhinv(\tau)\right|d\tau
  \label{eq:bound-cvar}
\end{equation}
となり，$[\alpha,1]$ における近似誤差はCVaRの誤差に一切寄与しない．
逆に一様重み $\Wone$ を最小化する FQF は，測度の上で大半を占める $[\alpha,1]$（長さ $1-\alpha$ の領域）の近似に多くの表現リソースを費やしていることになる（ただし $W_1$ の被積分関数の大きさは裾で大きくなるため，実際の誤差割合そのものが $1-\alpha$ となるわけではない点に注意を要する）．
また，Kiesel ら \cite{kiesel2016wasserstein} が指摘するように，CVaR（Expected Shortfall）は $W_1$ 距離に関して最悪 $1/\alpha$ の Lipschitz 定数をもつ．
したがって，同じ $\Wone(Z,\hat Z)$ の値であっても誤差が $[0,\alpha]$ に偏っていれば CVaR 誤差は最大 $1/\alpha$ 倍に拡大されうる．
これが，一様重みでの近似がCVaR推定に対して精度悪化を招きうる定量的根拠である．
\end{remark}
\subsection{離散近似への適用と正則化を含まない $g_\alpha$ による評価}
式\eqref{eq:zhat}の離散分布 $\hat Z_N$ に対する歪みリスク尺度は
\begin{equation}
  \rho_g(\hat Z_N)=\sum_{i=0}^{N-1}\left(g(\tau_{i+1})-g(\tau_i)\right)\Finv(\hat\tau_i)
  \label{eq:rho-discrete}
\end{equation}
と閉形式で計算できる．$\hat Z_N$ の分位点関数が階段関数であることから直ちに従う．
本研究では最適化時に後述の正則化歪み関数 $g_{\alpha,\lambda}$ を用いるが，
評価では常に式\eqref{eq:rho-discrete}に正則化を含まないCVaR歪み関数 $g_\alpha$（以下，純粋な $g_\alpha$ と呼ぶ）を代入して $\CVaR_\alpha(\hat Z_N)$ を計算する．
これにより，評価対象とは異なる目的関数を最適化したことによる見かけ上の性能向上という懸念を排除している．

%==================== 4 提案モデル ====================
\section{提案モデル：歪み重み付きquantile fraction提案}
\label{sec:proposed}

\subsection{FPN の損失の置換}
本研究の提案方針は極めて簡潔である．FQF の FPN が最小化する目的関数
$\Wone(Z,\hat Z_N)$ を，命題\ref{prop:bound}の右辺
\begin{equation}
  L_g(\tau)=\int_0^1\left|\Finv(u)-\Fhinv(u)\right|g'(u)\,du=\Wone(\dist Z,\dist\hat Z_N)
  \label{eq:Lg}
\end{equation}
に置き換える．ここで $\tau=(\tau_1,\dots,\tau_{N-1})$ はquantile fraction，$\hat Z_N$ は式\eqref{eq:zhat}であるが，
原子の位置 $\hat\tau_i$ は次項で述べるように中点から \textbf{$g$-中央値}（$g'$-重み付き中央値）へ変更する．
命題\ref{prop:bound}により，$L_g$ を小さくすることは $\rho_g$ の推定誤差の上界を小さくすることに直結する．
一様重み（$g(\tau)=\tau$）のときは $L_g=\Wone(Z,\hat Z_N)$ となり FQF の母集団目的に帰着するため，提案手法の目的関数は従来の FQF を特別な場合として包含している．

\subsection{最適quantile fraction＝歪み分布の1次元最適量子化}
\begin{proposition}[$g$-中央値原子]
\label{prop:gmedian}
quantile fraction $\tau_0<\dots<\tau_N$ を固定する．区間 $i$ の $L_g$ への寄与
$\ell^g_i(q)=\int_{\tau_i}^{\tau_{i+1}}|\Finv(u)-q|\,g'(u)\,du$ を最小化する原子は
\begin{equation}
  q_i=\Finv(\hat\tau_i),\qquad \hat\tau_i=g^{-1}\!\left(\frac{g(\tau_i)+g(\tau_{i+1})}{2}\right)
  \label{eq:gmedian}
\end{equation}
で与えられる．すなわち $\hat\tau_i$ は区間上の測度 $g'(u)du$ の中央値（$g$-中央値）である．
\end{proposition}
\begin{proof}
$\Finv$ は非減少であるから，$q=\Finv(m)$，$m\in[\tau_i,\tau_{i+1}]$ と書ける．
$\partial\ell_i^g/\partial q=\int_{\tau_i}^{\tau_{i+1}}\mathrm{sign}(q-\Finv(u))g'(u)du
=\left(g(m)-g(\tau_i)\right)-\left(g(\tau_{i+1})-g(m)\right)$ であり，これが零となるのは
$g(m)=(g(\tau_i)+g(\tau_{i+1}))/2$ のときである．$\ell^g_i$ は $q$ について凸なので，これが最小点である．
\end{proof}
一様 $g$ では式\eqref{eq:gmedian}は算術中点に一致し，FQF の補題1を特別な場合として含む．
なお，$\Finv$ が平坦な区間をもつ場合は $q_i$ は一意に定まらず，また $g$ が平坦な区間をもつ場合（$g_\alpha$ の $(\alpha,1]$ 部分）には $g$-中央値 $\hat\tau_i$ は一意に定まらない．式\eqref{eq:gmedian}における一般化逆関数による定義はその中で最小の代表元を選択する規約である．一方，最適化に用いる正則化歪み $g_{\alpha,\lambda}$（至る所 $g'>0$）では $\hat\tau_i$ は一意に定まる．

この命題と第\ref{subsec:bound}節の観察を合わせると，提案法の最適化問題は明確な意味をもつ．
$v=g(\tau)$ と変数変換すれば，$\hat Z_N$ の歪み分布 $\dist\hat Z_N$ は原子 $\Finv(g^{-1}(\hat v_i))$，重み $v_{i+1}-v_i$ の離散分布であり，
$L_g=\Wone(\dist Z,\dist\hat Z_N)$ は「歪み分布 $\dist Z$ を $N$ 点の離散分布で $L^1$ の意味で最良近似する」問題の目的関数である．
これは Graf and Luschgy \cite{graf2000quantization} の意味での確率分布の \textbf{1次元 $N$ 点最適量子化（$r=1$. $r$ は近似誤差の $L^r$ ノルムにおける次数）} にほかならない．
真の分位点関数 $\Finv$ が狭義単調（すなわち $Z$ に確率質量をもつ原子が存在しない）という条件下では，1次元において最適量子化器のVoronoi領域（各原子に最も近い点の集合）は区間となるため，分位点空間 $[0,1]$ の分割を区間 $[\tau_i,\tau_{i+1}]$ とする本研究のパラメータ化に一般性の損失はない．
ただし，この量子化同値が厳密に成り立つのは，真の分位点関数が固定された母集団・静的な最適量子化問題においてのみであり，深層ネットワークによる近似誤差やベルマン更新・方策変化に伴う分布変動を含む強化学習の動的設定全体に対して最適量子化が保証されるわけではない点に留意されたい．
最適量子化理論からは次の事実が得られる．
\begin{itemize}
  \item 最適量子化器は「原子＝各区間の中央値（$L^1$ の場合）」「境界＝隣接原子の中点」という Lloyd 型の固定点条件を満たす．
  \item $\dist Z$ に適切なモーメント条件および正則性条件（絶対連続性，有限の高次モーメント，密度の有界性など）が成り立つ場合，
  1次元の $N$ 点最適量子化誤差は $O(1/N)$ のオーダーをもつ（Zador 型定理）．
  有限1次モーメントのみから一般に $O(1/N)$ が従うとは限らないため，
  本研究で扱う各分布が既存の量子化定理の仮定を満たすかは別途確認を要する（第\ref{sec:future}節）．
\end{itemize}
上記の量子化誤差評価に必要な条件が成り立つ場合，命題\ref{prop:bound}から $\CVaR_\alpha$ の推定誤差にも $O(1/N)$ の上界が得られる．
CVaR歪み $g_\alpha$ の場合，$\dist Z$ は $Z$ を $Z\le\Finv(\alpha)$ に条件付けた分布であり，
その最適量子化は $N$ 点すべてを下側裾 $[0,\alpha]$ に配置する．
一様配置が同じ領域に置く点数はおよそ $\alpha N$ であるから，
両者の誤差の比はおおむね「点数の比」$1/\alpha$ になると期待される．
\begin{conjecture}[誤差上界と下界の比による $1/\alpha$ 改善]
\label{conj:alpha}
$\Finv$ が $(0,\alpha]$ で $C^1$ 級かつ $\Finv\in L^1[0,\alpha]$（$\CVaR_\alpha$ が有限）であり，さらに $(\Finv)'(\tau)=\tau^{-\beta}\ell(\tau)$（$\beta<2$，$\ell$ は緩変動関数．Student-$t_\nu$ では $\beta=1+1/\nu$，ガウス裾では $\beta=1$）の形で $\tau\downarrow0$ において正則変動する（regular variation）正則な分布族に対し，$\alpha N\ge2$ の領域において，一様quantile fractionによる式\eqref{eq:bound-cvar}の上界の下界と，歪み最適quantile fractionによる上界の比は，オーダーで $\Theta(1/\alpha)$ となる．
\end{conjecture}
なお，可積分性の仮定を導関数側（「$(\Finv)'$ が有界」「$(\Finv)'$ が $L^1[0,\alpha]$ に属する」等）に置くと，下側サポート非有界で $(\Finv)'(\tau)\to\infty$（$\tau\downarrow0$）となる分布（D1〜D3 はすべて該当）が除外されるため，可積分性は $\Finv$ 自身（$\CVaR_\alpha$ の有限性）に置き，導関数側は正則変動指数 $\beta<2$ で緩めている（$\beta<2$ は $\CVaR_\alpha$ の有限性と整合する）．
また，予想\ref{conj:alpha}は重み付き $\Wone$ 型の誤差\textbf{上界}に関する比である一方，注意\ref{rem:midpoint}で扱うのは CVaR の符号付き誤差そのものであり，両者は別の量である（後者の実測では $N$ 依存の改善比として観測される．注意\ref{rem:midpoint}参照）．
\begin{remark}[裾特異性と適応配置による二次精度]
\label{rem:midpoint}
正則化歪み $g_{\alpha,\lambda}$ は $[0,\alpha)$ 上で線形であるため，$g$-中央値は算術中点と一致し，歪み最適配置による CVaR 推定は $[0,\alpha]$ 上の\textbf{非一様格子}中点則に一致する．
本実験の分布はいずれも下側サポートが非有界で，$\Finv(\tau)\to-\infty$（$\tau\downarrow0$）の特異性をもつ（Student-$t$（自由度 $\nu$）では $\Finv(\tau)\sim -c\,\tau^{-1/\nu}$）．
\textbf{一様配置}では，この特異性を含む第1区間 $[0,1/N]$ の中点則誤差が
\begin{equation*}
  \int_0^h\tau^{-1/\nu}\,d\tau\;-\;h\,(h/2)^{-1/\nu}\;\propto\; h^{\,1-1/\nu}
\end{equation*}
の形となるため，CVaR 誤差は $O(N^{-(1-1/\nu)})$ に律速されると考えられる（$\nu=3$ で $O(N^{-2/3})$）．
実際，D3 では第1区間のみの寄与で M0 の全誤差の約 83〜95\%（$N\ge16$）を説明でき，その log--log 傾き $-0.68$ は理論値 $-(1-1/\nu)=-2/3$ と一致する．実測の D3 の M0 傾き $-0.72$ はこれと整合し，D1（混合ガウスの対数的な特異性）の $-1.27$ は特異性が弱いことを反映する．
一方，\textbf{歪み最適配置}は $L^1$ 最適量子化の適応性により，傾きの急な領域（裾）ほど区間幅を狭くとる．実測でも D3 の歪み最適配置の第1区間幅は $N=32$〜$256$ で $\tau_1\approx0.4\,N^{-3}$ と $N^{-\nu}$ で減少することを確認しており（一様の $1/N$ と比べて一挙に小さい），第1区間の特異性の寄与は $w_1^{1-1/\nu}\propto N^{-(\nu-1)}$ 程度まで抑えられる（$\nu=3$ で $O(N^{-2})$）．正則部分の中点則誤差 $O(N^{-2})$ と合わせ，提案法の符号付き CVaR 誤差は $O(N^{-2})$ に整合する減衰を示すと解釈できる（実測傾き $-1.79$〜$-1.84$．一般的な漸近次数の証明は今後の課題である）．
したがって改善比 $R$ は $N$ とともに増大する $N$ 依存の量であり，$1/\alpha^2$ は固定 $N$ の理論値ではなく漸近オーダーの目安である．
\end{remark}
\subsection{quantile fractionに関する勾配の導出：$g'(\tau_i)$ 倍}
\begin{proposition}[母集団目的における境界勾配]
\label{prop:grad}
$\Finv$ は連続かつ非減少であり，境界 $\tau_i$ および原子 $q_i=\Finv(\hat\tau_i)$ が式\eqref{eq:gmedian}の $g$-中央値に置かれている（母集団設定）とする．さらに各 $\tau_i$ は $g_{\alpha,\lambda}$ の微分可能点にある（すなわち $\tau_i\ne\alpha$）とする．このとき $i=1,\dots,N-1$ について
\begin{equation}
  \frac{\partial L_g}{\partial\tau_i}=g'(\tau_i)\left[2\Finv(\tau_i)-\Finv(\hat\tau_i)-\Finv(\hat\tau_{i-1})\right].
  \label{eq:dist-grad}
\end{equation}
\end{proposition}
\begin{proof}
$L_g=\sum_j\ell_j^g(q_j)$ において $\tau_i$ に依存するのは区間 $i-1$ と $i$ である．
FQF の場合（第\ref{subsec:fqf}節）と同様に，(a) 積分区間の端点としての依存と (b) 原子 $q_{i-1},q_i$ を通じた依存に分ける．
(b) は命題\ref{prop:gmedian}により $\partial\ell^g_j/\partial q_j=0$ であるから消える（包絡線定理：最適化後の目的関数をパラメータで微分するとき，最小点自身のパラメータ依存は寄与しない）．
(a) は Leibniz 則により，区間 $i-1$ の上端から $+|\Finv(\tau_i)-q_{i-1}|\,g'(\tau_i)$，区間 $i$ の下端から $-|\Finv(\tau_i)-q_i|\,g'(\tau_i)$ を与える．
$\hat\tau_{i-1}\le\tau_i\le\hat\tau_i$ と $\Finv$ の単調性により絶対値を外すと
\begin{equation*}
  g'(\tau_i)\left[\left(\Finv(\tau_i)-\Finv(\hat\tau_{i-1})\right)-\left(\Finv(\hat\tau_i)-\Finv(\tau_i)\right)\right]
\end{equation*}
となり，式\eqref{eq:dist-grad}を得る．
\end{proof}
すなわち，\textbf{母集団目的における提案法の境界勾配は FQF のquantile fractionに関する勾配（式\eqref{eq:fqf-grad}）に $g'(\tau_i)$ を乗じたものに一致する}．
実装では，真の分位点関数 $\Finv$ を分位点値ネットワークの出力 $\Finv_{\omega_2}$ で置き換え，
その出力への勾配を停止したplug-in型のFPN更新として式\eqref{eq:dist-grad}を用いる（全体の手順は\ref{subsec:impl}節）．
本命題は母集団設定における包絡勾配を導出したものであり，深層強化学習の実装においては，分位点値ネットワークに plug-in した擬似勾配更新（stop-gradient 型）として用いる．
学習途中のネットワーク出力が各区間の損失 $\ell_j^g$ の厳密な最小点である保証はなく，また FQF 原論文自身，深層ニューラルネットワークとベルマン更新を組み合わせた系の収束性は理論的に保証していない点に留意されたい．
また $\Finv$ や $g$ に平坦部がある場合，または $\tau_i=\alpha$（$g'_{\alpha,\lambda}$ の不連続点）にある場合には式\eqref{eq:dist-grad}は片側微分または劣勾配として理解する．
命題\ref{prop:grad}の適用条件をまとめると，(a) $g$ が非減少・正規化・絶対連続であること，(b) 考える境界で $g'(\tau_i)$ が存在すること（折点では劣勾配），(c) $\Finv$ が連続かつ非減少であること，(d) 原子が各区間の $g$-中央値に置かれていること，の4点であり，実装では至る所 $g'>0$ で区分線形な正則化歪み $g_{\alpha,\lambda}$ を用いるため (b) は $\tau_i\ne\alpha$ の管理で充足される．
勾配が零となる条件「$\Finv(\tau_i)$ が隣接原子の中点」は FQF と同じであるが，
原子の定義が変わっているため停留点の位置は異なる．
なお，$g'(\tau_i)>0$ である限り勾配に乗じられる係数だけを見れば零点は変化しないが，
提案法では代表点 $\hat\tau_i$ が算術中点から $g$-中央値へ変化するため勾配式の括弧内も変化し，
最終的な停留点配置は通常の FQF とは異なる．
学習ダイナミクスとしては，すでに $\tau_i<\alpha$ に位置する境界は大きな勾配係数 $(1-\lambda)/\alpha+\lambda$ を受ける一方，
$\tau_i>\alpha$ にある境界では係数が $\lambda$ にとどまるため，
初期配置によっては裾への移動が遅くなる可能性がある（実験Bの seed 依存性はこの現象と整合する）．
また，原子を $g$-中央値にせず算術中点のままにすると (b) の項が消えず，式\eqref{eq:dist-grad}は厳密な勾配とならずその近似にとどまる．

\subsection{正則化歪み関数}
CVaR 歪み $g_\alpha$ は $(\alpha,1]$ で $g'_\alpha=0$ であるため，そこにあるquantile fractionの勾配が零となって動かなくなり，
また $g_\alpha^{-1}$ が $[\alpha,1]$ で定義できず，$g$-中央値が定義できなくなる．
これを避けるため，最適化には正則化歪み関数
\begin{equation}
  g_{\alpha,\lambda}(\tau)=(1-\lambda)\,g_\alpha(\tau)+\lambda\tau,\qquad
  g'_{\alpha,\lambda}(\tau)=\begin{cases}(1-\lambda)/\alpha+\lambda & \tau<\alpha\\ \lambda & \tau>\alpha\end{cases}
  \label{eq:g-reg}
\end{equation}
を用いる（既定 $\lambda=0.05$）．至る所 $g'\ge\lambda>0$ であり，$\rho_{g_{\alpha,\lambda}}=(1-\lambda)\CVaR_\alpha+\lambda\E$ は
CVaR と期待値の凸結合である．評価は前述のとおり純粋な $g_\alpha$ で行う．
$\lambda$ は「分布全体をどの程度保持するか」を制御するパラメータでもあり，RL への適用では裾以外の情報が方策改善に必要になりうるため重要な役割を果たす．

\subsection{FQF への実装（歪み FQF）}
\label{subsec:impl}
以上をまとめると，提案法（以下，歪み FQF と呼ぶ）の1更新は以下の手順で与えられる．
\begin{enumerate}
  \item 状態特徴 $\psi(s)$ から FPN によりquantile fraction $\tau_0=0<\tau_1<\dots<\tau_N=1$ を得る（ソフトマックス＋累積和）．
  \item 各区間の $g$-中央値 $\hat\tau_i=g_{\alpha,\lambda}^{-1}\!\left((g_{\alpha,\lambda}(\tau_i)+g_{\alpha,\lambda}(\tau_{i+1}))/2\right)$ を計算する．
  \item 分位点値ネットワークで $\Finv_{\omega_2}(\hat\tau_i;s,a)$ を出力し，目標分布との Huber 分位点回帰損失（式\eqref{eq:qr-loss}）で $\omega_2$ を更新する．
  \item 式\eqref{eq:dist-grad}の plug-in 型勾配（$\Finv$ を分位点値ネットワークの出力で置き換え，その出力への勾配を停止したもの）により $\partial L_g/\partial\tau_i$ を計算し，ソフトマックス・累積和のヤコビアンで連鎖して FPN の $\omega_1$ を更新する（分位点値ネットワークには勾配を流さない）．
  \item 行動選択は $a^*=\arg\max_a\rho_{g_\alpha}(\hat Z(s,a))$（式\eqref{eq:rho-discrete}，純粋 $g_\alpha$）で行う．
\end{enumerate}
手順2と4が FQF からの変更点であり，追加計算量は $O(N)$ である．

\subsection{新規性調査の結果要約}
本提案の新規性について，(a) FQF とリスク尺度の結合，(b) $\tau$ のサンプリング・サポート点配置の $\Wone$ 最適化，
(c) 歪み重み付き $\Wone$ と損失設計，(d) 統計・最適輸送における歪み尺度の最適量子化，の4観点で文献調査を行った．
調査方法は，arXiv・Google Scholar を主要データベースとし，上記4観点に対応する検索語群（distributional reinforcement learning，risk measure，distortion，quantile optimization，Wasserstein，CVaR，spectral risk，quantile fraction など）を組合せて 2026 年 9 月時点までの文献を検索・精査し，約 70 文献を関連度ごとに 10 グループに分類したものである（inclusion 基準は「分布型 RL の表現設計・リスク尺度・1次元量子化のいずれかに直接関与」）．
結論として，\textbf{本稿で実施した文献検索の範囲では，FPN の目的関数そのものを歪み重み付き Wasserstein 距離に置換し，quantile fraction学習を歪み分布の最適量子化として定式化した先行研究は確認できなかった}．
近接する先行研究との差分を表\ref{tab:novelty}に示す．
Wasserstein 距離に基づく確率分布の最適離散化は Pflug and Pichler \cite{pflug2014multistage} などの応用確率論の系譜で広く研究されてきたが，これを分布型強化学習における fraction 学習機構と直接結びつけた研究は見当たらない．
また，近年 Gan ら \cite{gan2026energy} は FQF とスペクトルリスク尺度（SRM）を組み合わせたエネルギー管理手法（FQF-SRM）を提案しているが，これは学習された FQF の出力から SRM に基づいて行動選択を行うものであり，FPN の目的関数は依然として一様重み $W_1$ のままである．
本研究の差分は，リスク尺度を行動選択時だけでなく \textbf{FPN の近似損失そのものに導入し}，fraction 配置を歪み分布の最適量子化として導出・学習させる点にある．
金融工学分野では，深層強化学習を用いたデリバティブのヘッジング（Cao et al. \cite{cao2021deep,cao2023gamma}）や，裾の極値分布（GPD）を用いて極端分位点の予測精度を高める EX-DRL（Malekzadeh et al. \cite{malekzadeh2024exdrl}）などが提案されているが，いずれも分位点回帰のサンプリングやネットワーク構造の工夫であり，fraction の最適量子化による解像度再配分とは異なる．
なお，本調査の結論は検索範囲に依存した限定的な知見であることに留意されたい．
補完確認（2026-09-12）では，quantile geometry の正則化と Wasserstein ロバストな分位推定を扱う RQIQN（Zhang ら \cite{zhang2026rqiqn}）が把握されたが，同手法は Bellman 標的への fraction 依存の補正による推定の頑健化であり，目的関数をリスク尺度に応じて置換して fraction 配置を最適化する本提案とは目的・機構が異なる．

\begin{table}[htbp]
\centering
\caption{近接する先行研究との差分}
\label{tab:novelty}
\small
\setlength{\tabcolsep}{3pt}
\begin{tabular}{p{3.8cm}p{5.0cm}p{5.6cm}}
\toprule
文献 & 関連する内容 & 本提案との差分\\
\midrule
FQF \cite{yang2019fqf} & quantile fractionを一様 $\Wone$ 最小化で学習 & 目的関数がリスク尺度と無関係．本提案は $g$ で重み付け\\
FQF-SRM \cite{gan2026energy} & FQF とスペクトルリスク尺度（SRM）による行動選択 & FPN の損失は一様 $W_1$ のまま．本提案は FPN 損失そのものを置換\\
IQN \cite{dabney2018iqn} & 行動選択時に $\tau$ のサンプリング分布を歪めてリスク感応方策 & 学習時の $\tau$ 配置は一様．表現の解像度は変えない\\
Beyond CVaR \cite{moghimi2025beyond} & 静的スペクトルリスク尺度を最適化する分布型RL & 最適化対象は方策．近似ノードの再配置は行わない\\
Tail-Safe \cite{zhang2025tailsafe} & ヘッジングで $\tau$ サンプリングを裾に集中させる TCC & 発見的な温度制御であり，Wasserstein 最適量子化に基づく理論的特徴づけをもたない\\
EX-DRL \cite{malekzadeh2024exdrl} & GPD を用いた極端分位点のモデル化によるヘッジング & 極値理論による裾モデリングであり，quantile fraction の最適量子化ではない\\
最適離散化理論 \cite{graf2000quantization,pflug2014multistage} & Wasserstein 距離に基づく確率分布の最適量子化 & 静的な確率論の枠組み．RL のquantile fraction学習へ接続した点に本提案の独自性がある\\
Lim and Malik \cite{lim2022distributional} & リスク感応方策のための分布型RL（CVaR の時間不整合の指摘を含む） & 近似表現の設計ではなく方策最適化の枠組み\\
RQIQN（Zhang ら, 2026）\cite{zhang2026rqiqn} & quantile geometry の正則化（Wasserstein ロバストな分位推定・fraction 依存補正） & 目的関数は変更しない頑健化手法．本提案は FPN 損失そのものを歪み重み付きに置換\\
Rockafellar and Uryasev \cite{rockafellar2000cvar} & CVaR の変分表現と凸最適化 & 分布近似誤差の議論は含まない\\
\bottomrule
\end{tabular}
\end{table}
%==================== 5 数値実験 ====================
\section{数値実験}
\label{sec:exp}
計算環境はすべて AMD Ryzen 7 5825U（8コア16スレッド），32\,GB RAM，CPU演算のみ（PyTorch CPU版）である．

\subsection{実験A：歪み分布の1次元最適量子化}
\label{subsec:expA}
実験Aは，真の分布が既知で，その分位点関数 $\Finv$ と部分期待値が高精度に評価できるときに，quantile fraction配置の違いだけで $\CVaR_\alpha$ の推定誤差がどれだけ変わるかを測る．
学習の不確定性を排し，提案法の原理（第\ref{sec:proposed}節）そのものを検証する目的である．
この意味で実験Aは，M2 が真の $\Finv$ と閉形式部分期待値を用いる \textbf{真値参照の歪み最適量子化実験}であり，深層ネットワークとベルマン更新を含む学習（実験B）とは性質の異なる検証であることに留意されたい．

\subsubsection{実験手順}
\begin{enumerate}
  \item \textbf{対象分布}（いずれも損益であり，下側裾がヘビー）を用意する．
  D1：ガウス混合 $0.9\,\mathcal N(0,1)+0.1\,\mathcal N(-4,1)$，
  D2：負の対数正規 $Z=e^{\sigma^2/2}-\exp(\sigma\xi)$，$\xi\sim\mathcal N(0,1)$，$\sigma=1$（平均0），
  D3：Student-$t$（自由度 $\nu=3$）．
  \item \textbf{参照真値}を高精度に計算する．部分期待値 $\mathrm{PE}(a,b)=\E[Z\,\mathbb{I}\{a<Z\le b\}]$ の閉形式（ガウス成分・$t$分布・対数正規）と，分布関数の数値反転（許容誤差 $10^{-12}$ の数値求根）による分位点 $\Finv(\alpha)$ を用い，
  $\CVaR_\alpha(Z)=\alpha^{-1}\mathrm{PE}(-\infty,\Finv(\alpha))$ とする（ガウス混合の $\Finv$ 自身は単純な閉形式をもたない）．$10^7$ サンプルのモンテカルロ（seed 0）で $|\text{参照値}-\text{MC}|<3\,\mathrm{SE}$（SE は MC 推定量の標準誤差）をクロスチェックする．
  \item \textbf{格子}を $N\in\{8,16,32,64\}$（任意で $128,256$），$\alpha\in\{0.01,0.05,0.1\}$ とする．
  \item \textbf{比較手法}として以下を実装する．
  M0：一様配置 $\tau_i=i/N$，算術中点原子（FQF相当）．
  M1：$g$-一様配置（$g$-uniform）$\tau_i=g^{-1}(i/N)$（裾の事前知識のみを与えた対照，最適化なし）．
  M2：数値最適化で得た候補解（以下，便宜上「歪み最適配置」と呼ぶ．大域最適性は保証しない）．$\tau=\mathrm{cumsum}(\mathrm{softmax}(\theta))$（$\theta$ は最適化変数であるソフトマックス前の実数ベクトルで，QR-DQN の原子 $\theta_i$ とは別物）とパラメータ化し，
  式\eqref{eq:dist-grad}の解析勾配をソフトマックスのヤコビアンで連鎖して L-BFGS-B（制約付き準ニュートン法ソルバ，maxiter 500，gtol $10^{-10}$）で $L_g$ を最小化する．
  初期値は M1 由来の $\theta$ と Dirichlet 摂動（非負で総和1に正規化した乱数ベクトル）4つの計5スタートとし，$L_g$ 最小の解を採用する．
  M3：Lloyd 型固定点（原子 $\leftarrow$ 区間の $g$-中央値，境界 $\leftarrow$ 隣接原子の中点，200反復）を M2 のクロスチェックに用いる．
  \item \textbf{目的関数}は台形則を使わず厳密に評価する．CVaR 歪みでは $g'$ が区分定数なので $\alpha$ を跨ぐ区間 $[\tau_i,\tau_{i+1}]$ を $\alpha$ で分割し，各小区間の積分を部分期待値で閉形式計算する．
  最適化には正則化歪み $g_{\alpha,\lambda}$（$\lambda=0.05$）を用いる．
  \item \textbf{評価}は純粋な $g_\alpha$ で行う．誤差 $e=|\CVaR_\alpha(\hat Z_N)-\CVaR_\alpha(Z)|$，改善比 $R=e_{\mathrm{M0}}/e_{\mathrm{M2}}$，$[0,\alpha]$ 内のquantile fraction数を記録する．
  \item \textbf{点数換算値（N-equivalence）}：M2 の $N$ 点と同等以下の誤差を一様配置が達成するのに必要な最小点数 $N'$ を求める．手順は，まず $2^n$（$n\le10$）の格子上で M2 の誤差を初めて下回る最小の格点で区間を特定し（bracket），その隣接格点の間の整数点数に対して二分探索により最小の $N'$ を求める．
  \item \textbf{収束レート}：$\alpha=0.05$ について，主格子（$N\in\{8,\dots,256\}$）と拡張測定（$N\in\{256,512,1024\}$，$N=256$ は主格子の値を優先）を併せた $N=8$〜$1024$ のデータで，$\log e$ を目的変数，$\log N$ を説明変数とする最小二乗回帰（OLS）の傾きと決定係数 $R^2$ を分布ごとに求める．
  \item \textbf{有限サンプル設定 A-fs}：$\Finv$ を，$n\in\{10^3,10^4,10^5\}$ サンプルから作る経験分布（各サンプルに質量 $1/n$ を置いた分布）の分位点関数（経験分位関数）で置き換えて（seed 0--9 で試行），M2 を経験 $\Finv$ 上で最適化し，参照真値と比較する．
  \item \textbf{アブレーション}（構成要素の一部を変えて効果を確認する追加実験）：D1 で $\lambda\in\{0,0.01,0.2\}$，および Wang 変換 $g_\eta(\tau)=\Phi(\Phi^{-1}(\tau)+\eta)$（$\Phi$ は標準正規累積分布関数，$\eta=1.0$；なお $\eta>0$ は歪みの強さを表し，§\ref{subsec:cvar}の閾値変数 $\eta$ とは無関係）を試す．
\end{enumerate}

\subsubsection{結果}
図\ref{fig:placement}は D1，$N=16$，$\alpha=0.05$ におけるquantile fraction配置を分位点関数上に示したラグプロット（rug plot. 各データ点を軸上に短い縦線で示すプロット）である．
一様配置 M0 は裾 $[0,\alpha]$ に1点も置かないのに対し，提案法 M2 は15個の内部quantile fractionのうち13個を $[0,\alpha]$ 内に配置し，
残りを $\lambda$ による重みで分布の残部に割り当てている．これは第\ref{sec:proposed}節の「歪み分布の最適量子化」という特徴づけの直接の可視化である．

\begin{figure}[htbp]
\centering
\includegraphics[width=0.85\linewidth]{fig_placement.png}
\caption{D1（ガウス混合），$N=16$，$\alpha=0.05$ におけるquantile fraction配置のラグプロット．上から M0（一様），M1（$g$-一様配置），M2（歪み最適）．網掛けは下側裾 $[0,\alpha]$．}
\label{fig:placement}
\end{figure}

表\ref{tab:A1}に $\alpha=0.05$，$N=16$ での結果を示す．3分布すべてで改善比 $R$ は主結果表の判定基準として事前に定めた $R\ge2$ を大きく上回った．
なお，$N=16, \alpha=0.05$ の一様配置 M0 では最初の区間 $[0, 1/16]$ が $[0, \alpha]$ を完全に覆うため，CVaR 推定が単一分位点 $\Finv(1/32)$ \textbf{のみで行われる特殊設定}であることに留意されたい．
D1 で $R=97.9$ に達するのは，混合成分 $\mathcal N(-4,1)$ による損失クラスタが $\tau<0.1$ に存在し，
一様配置の単一分位点ではこの構造がほぼ捉えられないためである．
より客観的な比較として，後述の仮説 H1 対象領域 $\alpha N\ge2$ における最小改善比は $R=42.3$ であった．
対照 M1 に対しても提案法はすべての設定で上回り（$e_{\mathrm{M2}}/e_{\mathrm{M1}}=0.05$〜$0.17$，成功基準 $\le0.7$ を充足），
quantile fractionの最適化が裾への事前配置を超える効果をもつことを示した．
なお，Lloyd 型固定点 M3 によるクロスチェックでは，全 54 設定（3分布$\times$3水準$\times$6格子）のうち 46 設定で M2 の目的関数値 $L_g$ が M3 以下であり，
残る 8 設定でも相対差は最大約 22\% にとどまった．
このことから，M2 のマルチスタート L-BFGS-B がすべての設定で大域最適に達した保証はない（M2/M3 の $L_g$ の最小値を採用する等の追加確認は将来課題とする）．
また，参照真値は $10^7$ サンプルのモンテカルロ（seed 0）との照合により，3分布$\times$3水準の全 9 構成で
$|\text{参照値}-\text{MC}|<3\,\mathrm{SE}$ を満たすことを確認した（最大 $|z|=1.69$. $z=(\text{参照値}-\text{MC})/\mathrm{SE}$ は標準化偏差）．

\begin{table}[htbp]
\centering
\caption{実験A主結果（$\alpha=0.05$，$N=16$）．$e$ は $\CVaR_\alpha$ 推定の絶対誤差，$R=e_{\mathrm{M0}}/e_{\mathrm{M2}}$．M1 は裾の事前知識のみを与えた対照（$g$-一様配置，最適化なし）．M2 は真の $\Finv$ を用いる真値参照設定である．判定列 $R\ge2$ は全格子での最低基準である．}
\label{tab:A1}
\small
\setlength{\tabcolsep}{4pt}
\begin{tabular}{lcccccc}
\toprule
分布 & $e_{\mathrm{M0}}$ & $e_{\mathrm{M1}}$ & $e_{\mathrm{M2}}$ & $R$ & 判定（$R\ge2$） & M2の $[0,\alpha]$ 内quantile fraction数\\
\midrule
D1 ガウス混合 & $0.309$ & $0.0626$ & $0.00316$ & 97.9 & PASS & 13 / 15（M0 は 0）\\
D2 負の対数正規 & $2.12$ & $0.245$ & $0.0404$ & 52.4 & PASS & 13 / 15\\
D3 Student-$t$（自由度3） & $0.974$ & $0.148$ & $0.0246$ & 39.6 & PASS & 13 / 15\\
\bottomrule
\end{tabular}
\end{table}

図\ref{fig:error}は誤差 $e$ 対 $N$ の log-log プロットである．
仮説の判定基準は2階層である：全格子での最低改善基準 $R\ge2$（表\ref{tab:A1}），および主仮説 H1（安定領域 $\alpha N\ge2$ かつ $N\ge32$ の格子点，すなわち $\alpha\in\{0.05,0.1\}\times N\in\{32,64\}$ から $\alpha=0.05,\,N=32$（$\alpha N=1.6$）を除いた集合での $R\ge5$）である．
観測結果は，対象領域における最小の改善比が $R=42.3$ であり H1 を充足した（表\ref{tab:A1}の判定列は全格子での最低基準であり，H1 とは水準が異なる）．
収束レートは，$\alpha=0.05$，$N=8$〜$1024$ の併合データに対する経験的 log--log 傾きとして，歪み最適配置 M2 で $-1.79$〜$-1.84$（決定係数 $R^2\ge0.998$），一様配置 M0 で $-0.72$〜$-1.27$（同 $R^2\ge0.994$，対照 M1 も $-0.69$〜$-1.29$ と同様）であり，M2 はおおむね二次減衰に，M0 はおおむね一次減衰に整合した（注意\ref{rem:midpoint}）．
一般的な漸近次数の証明は今後の課題である．なお $N=1024$ における M2 の誤差は $1.8\times10^{-6}$〜$1.7\times10^{-5}$ であり，参照真値の計算精度（数値求根の許容誤差 $10^{-12}$）や倍精度浮動小数点の精度に比べて十分大きく，誤差測定が数値精度の下限に達していないことは確認済みである．
改善比 $R$ は $N$ とともに増大し（$\alpha=0.05$，$N=16$ で $39.6$〜$97.9$，$N=1024$ で $1.1\times10^3$〜$4.2\times10^3$），固定 $N$ での理論比 $1/\alpha^2=400$ との「一致」を主張するものではない．

\begin{figure}[htbp]
\centering
\includegraphics[width=0.95\linewidth]{fig_error.png}
\caption{$\CVaR_\alpha$ 推定誤差 $e$ 対 $N$（log-log）．列：D1/D2/D3，行：$\alpha\in\{0.01,0.05,0.1\}$，線：M0/M1/M2．破線は参照傾き $-1$．}
\label{fig:error}
\end{figure}

N-equivalence（図\ref{fig:nequiv}）では，$\alpha=0.05$ において，D1 ガウス混合では M2 の16点と同等の誤差を一様配置が達成するには約695点が必要であった（$N=8$ の場合は約257点）．
一方，D2・D3 では $N=16$ における M2 の達成した誤差が極めて小さく，一様配置は探索上限の1024点を用いても同等の精度に達しなかった（図では数値を省略）．
D1 における約695点という値は，「quantile fraction数を約43倍にする」のと「配置を変える」のが同等の効果をもつことを意味し，
一様配置側の減衰が下側裾の特異性により一次程度にとどまること（注意\ref{rem:midpoint}）を反映している．

\begin{figure}[htbp]
\centering
\includegraphics[width=0.95\linewidth]{fig_nequiv.png}
\caption{点数換算値（N-equivalence，$\alpha=0.05$）．M2 の $N$ 点と同等以下の誤差を一様配置 M0 が得るのに要する点数 $N'$（探索上限は1024点，上限に達した場合は数値を省略）．}
\label{fig:nequiv}
\end{figure}

有限サンプル設定 A-fs（$n=10^4$）では，10シード平均の改善比の最小値は $R=2.83$，M2 が M0 を上回ったシードの割合は 10 中 8 以上であった．
参照真値を用いる設定に比べて改善幅が縮むのは，経験分位関数の下側裾（$\tau<\alpha$ の領域はサンプルの5\%以下）が統計的に不安定であり，
M2 がその不安定な裾に多くの点を配置するため経験分位関数の裾の推定分散に対する感度が高まるからである．
それでも多くの設定で M0 を上回っており，実際の RL（分位点ネットワークが有限データから学習する）でも改善が期待できる根拠となる．
$\lambda$ アブレーション（D1，$N=16$，$\alpha=0.05$）では，$\lambda=0$（純粋 CVaR）で $R=129.6$ と最も改善が大きい一方，
一様重み $\Wone(Z,\hat Z_N)$ の値は $3.69$ と大きく悪化し，$\lambda=0.01$ で $R=114.6$（$\Wone=0.75$），
既定の $\lambda=0.05$ で $R=97.9$（$\Wone=0.46$），$\lambda=0.2$ で $R=67.2$（$\Wone=0.27$）となった
（$[0,\alpha]$ 内quantile fraction数は順に 15，14，13，11 個）．
ここで $\lambda=0$ では $g_\alpha^{-1}$ が $(\alpha,1]$ で一意に定まらず裾外の境界勾配も零となるが，本実装では $\tau=\mathrm{cumsum}(\mathrm{softmax}(\theta))$ のパラメータ化により全区間幅が厳密に正に保たれ，$g$-中央値は一般化逆関数の最小代表元規約，境界は $[\varepsilon,1-\varepsilon]$ クリッピング（$\varepsilon=10^{-6}$）で運用したため，最適化自体は数値的に進行し，裾外の境界は勾配が零のまま初期値近傍に留まる（すなわち $\lambda=0$ の配置は非同定的であり，同値類の中の1解が得られている）．
これは $\lambda$ が「裾の推定精度」と「分布全体の近似品質」のトレードオフを制御することを定量的に示す．
また Wang 変換（$\eta=1.0$，$N=16$）による歪みリスク尺度 $\rho_{g_\eta}$ の推定誤差は，
D1 で $3.1\times10^{-4}$（$R=483.5$），D2 で $0.039$（$R=29.9$），D3 で $0.020$（$R=28.3$）であり，
CVaR 以外の歪みリスク尺度に対しても提案法が同様に機能することを確認した．

\subsection{実験B：toy RL プロトタイプ}
\label{subsec:expB}
実験Bは，提案法を実際に FQF に実装し，学習された分布から計算した $\CVaR_\alpha(s,a)$ の推定精度を測る．

\subsubsection{環境 RiskyChain-v0}
自作環境 RiskyChain-v0（gymnasium API，離散行動）を用いる．
ホライズン（1エピソードあたりの最大ステップ数）$H=10$，位置 $x\in\{0,\dots,9\}$，$\gamma=1$ である．観測は $\mathrm{onehot}(x)\oplus t/H$ の11次元ベクトルである．ここで $\mathrm{onehot}(x)$ は位置 $x$ に対応する1成分だけが1で他が0の10次元ベクトル（例：$x=3$ なら $(0,0,0,1,0,0,0,0,0,0)$），$t$ は開始からの経過ステップ数，$\oplus$ は2つのベクトルを直列に結合する操作（連結，concatenation）を表す．10次元の $\mathrm{onehot}(x)$ の後ろに1次元の値 $t/H\in[0,1]$（経過時間の割合）を結合するため，観測は11次元となる．
行動 SAFE は決定的報酬 $0.1$ を与え位置は不変，行動 RISKY は確率 $1-p(x)$ で成功（報酬 $0.5$，$x\leftarrow x+1$），
確率 $p(x)=0.002(x+1)$ でクラッシュし，ヘビーテイル損失 $-5(1+|T_3|)$（$T_3$ は自由度3の Student-$t$ 分布に従う確率変数）を受けて位置は不変となる．
到達可能な状態は $55$ 個であり，すべての $x$ で $\E[\mathrm{RISKY}]>0.1$ のためリスク中立最適方策は常に RISKY である一方，
常に RISKY を選ぶとクラッシュが1回以上起こる確率は約 $10.4\%>\alpha=5\%$ となる．
したがって $\CVaR_{0.05}$ を目的とした場合，「最初の $k$ ステップのみ RISKY，以降 SAFE」とする11通りの閾値方策のクラス内で評価した結果，
クラッシュを一切起こさず確定的に収益 $1.0$ を得る $k^*=0$（all-SAFE）がクラス内最良方策となり，その CVaR は確定値 $1.0000$（誤差ゼロ）となる（なお，$k=1$ の解析値は $0.9595<1.0$ であり，先行の有限モンテカルロ試行でのばらつきはサンプル分散 $\mathrm{SD}\approx0.16$ による推定ノイズであったことが判明している）．
これは一般の状態依存方策や履歴依存方策を含む全方策クラスに対する大域的最適性を意味しない．
以降，この方策を閾値方策クラス内最良（$k^*=0$）と呼ぶ．
リスク中立とリスク回避で最適方策が乖離する点が，この環境の設計意図である．
なお，本環境の収益分布（特に all-SAFE 方策下では確定値 $1.0$ の点質量）をもつため，真の分位点関数は狭義単調でなく，第\ref{sec:proposed}節の量子化同値や命題の前提（狭義単調性）は厳密には成立しない．$g$-中央値の非一意性により FPN の停留点が一意でなくなりうる点は，実験Bの seed 依存性の一因となりうる（式\eqref{eq:gmedian}の一般化逆関数規約，および命題\ref{prop:gmedian}直後の注意を参照）．

\subsubsection{実験手順}
\begin{enumerate}
  \item \textbf{共通ネットワーク構成}：MLP $11\to256\to256\to N\cdot|\mathcal A|$（ReLU），$N=16$，Adam（lr $10^{-4}$，FPN lr $2.5\times10^{-5}$），
  Huber 分位点回帰損失（$\kappa=1$），リプレイバッファ容量 50k，バッチ 64，ターゲット網のハード更新 500 ステップごと，
  $\varepsilon$-greedy で $\varepsilon$ を $1.0\to0.05$ に減衰（20k ステップで線形），
  予算は 100k 環境ステップ，seed $\{0,1,2\}$．
  なお，FQF 原論文における FPN 最適化のエントロピー正則化や RMSprop は本実験では用いず，Adam による標準的最適化としている．
  ベルマン標的に用いる次状態の行動は，ターゲットネットワークが推定する $\CVaR_{0.05}$ を最大化する行動（CVaR-greedy）として選択した．
  \item \textbf{Bellman 標的の原子重み付け}：ベルマン標的を構成する際，FQF 系（B0/B1/B1g）はターゲット側の離散分布の原子を区間幅で重み付けして損失を計算する（FPN が定義する離散分布を忠実に反映）．
  すなわち，ターゲット側の境界 $\tau'_j$，原子 $q'_j$（$j=1,\dots,N$）と現在側の level $\tau_i$・原子 $q_i$ に対し，
  \begin{equation}
    \mathcal{L}_{\mathrm{QR}}=\frac1N\sum_{i=1}^N\sum_{j=1}^N w'_j\,\rho^\kappa_{\tau_i}\!\left(r+\gamma q'_j-q_i\right),\qquad w'_j=\tau'_{j+1}-\tau'_j
    \label{eq:b-loss}
  \end{equation}
  を最小化する（$\rho^\kappa_\tau$ は式\eqref{eq:huber-quantile}の分位点 Huber 損失）．
  一方，IQN 系（B2/B2p）はターゲット $\tau'_j$ を学習時のサンプリング分布（B2 は一様，B2p は歪み分布）からサンプルし，等重み $w'_j=1/N$ で平均する（重要度重みは付けない）．
  \item \textbf{比較手法}：
  \begin{itemize}
    \item B0 FQF：一様 $\Wone$ のquantile fractionに関する勾配（一様初期化）．
    \item B1 歪み FQF（一様初期化）：提案法の境界勾配（式\eqref{eq:dist-grad}）と $g$-中央値原子を採用．初期化は一様配置．
    \item \textbf{B1g 歪み FQF（$g$-一様初期化・主結果）}：B1 と同一だが，FPN の初期 bias を $g$-一様配置 $\tau_i = g^{-1}_{\alpha,\lambda}(i/N)$（実験Aの M1 相当）に初期化し，初期から下側裾に解像度を配分する．
    \item B2 リスク感応 IQN：学習時 $\tau\sim U(0,1)$ 16点，コサイン埋め込み64次元，行動選択時 $\tau\sim U(0,\alpha)$．
    \item \textbf{B2p IQN（学習時歪み $\tau$）}：B2 と同一だが，\textbf{学習時のサンプリング（現在側の $\tau_i$ とターゲット側の $\tau'_j$ の双方）}のみ $\tau = g^{-1}_{\alpha,\lambda}(U(0,1))$ に変更し（分位点 Huber 損失の level には現在側 $\tau_i$ を，ターゲット原子の生成にはターゲット側 $\tau'_j$ を用い，重みは等重み $1/N$ のまま重要度重みは付けない），学習時から解像度を下側裾に集中させる対抗手法．行動選択（$\tau\sim U(0,\alpha)$）と決定的評価グリッドは B2 と完全同一である．
  \end{itemize}
  \item \textbf{方策（全手法共通）}：$\CVaR_\alpha$-greedy（$\alpha=0.05$）．すべての手法で行動選択および評価に同じリスク水準 $\alpha=0.05$ を用いる．
  \item \textbf{評価}：5k ステップごとに (i) \textbf{M-est（主指標）}：greedy 方策を固定し，55状態$\times$2行動のすべてについて $10^4$ 本のロールアウトに基づくモンテカルロ推定値 $\widetilde{\mathrm{CVaR}}_\alpha(s,a)$ を参照値として計算し，
  ネットワークの $\CVaR_\alpha$ との MAE を測る（閾値方策クラス内最良が切り替わる「意思決定に関わる状態」に限定した MAE も報告）．
  なお，方策のクラッシュ確率が異なるため，参照値の MC 標準誤差（SE）は手法間で異なる（実測: all-SAFE 系 SE 平均 0.035，all-RISKY 系 0.111，意思決定状態では 0.021 / 0.133）．手法間の大きな MAE 差に比べれば相対順序を覆すほどではないが，絶対値の解釈には注意を要する．
  さらに重要な交絡として，\textbf{推定対象となる収益分布そのものが手法間で異なる}：各手法の greedy 方策が異なるため（B0 は all-RISKY，B1g/B2 はほぼ all-SAFE），B0 は重い裾をもつ分布を，B1g/B2 はほぼ確定値に近い分布を推定している．この交絡を検討する共通方策下の交差評価を第\ref{subsec:cross}節に追加した．
  \item (ii) \textbf{M-pol}：5{,}000 評価エピソード（$\varepsilon=0$）の収益の経験 $\CVaR_{0.05}$ および平均収益，
  (iii) \textbf{quantile fraction診断}：100k ステップ時点での全評価状態にわたる学習済み $\tau$ の $[0,\alpha]$ 内のquantile fraction数．
  \item \textbf{IQN 系の CVaR 推定方法}：B0・B1・B1g は 16 個の重み付き原子から式\eqref{eq:rho-discrete}により CVaR を算出する．
  一方 B2・B2p（IQN）は固定原子分布をもたないため，評価には $\tau$ を $(0,\alpha)$ 内の等間隔16点 $\tau_j=\alpha(2j-1)/(2N)$ に固定した決定的グリッドを用い，
  出力分位点値の平均を CVaR 推定値とする．
  この読み出しは裾に16点すべてを割り当てられるため，B1g の「裾に約14個」という配置とは読み出し構造が異なり（IQN の連続 $\tau$ 入力がもつ本来の柔軟性である），B2 の有利性の一部はこの読み出し差に由来する可能性がある（第\ref{subsec:cross}節の交差評価とあわせて解釈されたい）．
\end{enumerate}

\begin{remark}[実験Bが測るもの]
状態ごとの $\CVaR$-greedy 制御は静的な $\CVaR$ 目的と時間整合しない（Lim and Malik \cite{lim2022distributional} も指摘する CVaR の時間不整合性）．
したがって実験Bが測るのは\textbf{分布の推定品質であり，方策の最適性ではない}．
さらに，M-est の参照値は各手法の greedy 方策の下で計算されるため，手法間の M-est 差には「表現方式の差」と「推定対象となる収益分布の難易度の差（B0 は重い裾，B1g/B2 はほぼ確定値）」が混在し，純粋な表現精度の比較ではない．この交絡の検討は第\ref{subsec:cross}節の共通方策交差評価で扱う（ただし本評価にも設計上の限界がある．同節参照）．M-pol は参考値として報告する．
\end{remark}

\subsubsection{結果}
表\ref{tab:B1}に100kステップ時点の最終指標を示す（3シード平均$\pm$標本SD，ddof=1）．
主結果である $g$-一様初期化入りの提案法（B1g）は，全110対の M-est MAE において $0.96\pm0.51$ を達成し，FQF（B0）の $7.02\pm0.14$ から約7分の1に誤差を低減した（図\ref{fig:B2}）．
意思決定に関わる状態に限定した MAE では B0 の $14.41\pm0.15$ に対し B1g は $0.39\pm0.29$ と約37分の1の大幅な改善を示した．
下側裾 $[0,\alpha]$ 内のquantile fraction数は状態あたり $14.06\pm0.11$ 個（15個中）に達し，3シードすべてで理論配置（約14個）を忠実に実現している（図\ref{fig:B3}）．
旧 B1（一様初期化）では seed 0 が裾に移動できず（裾0個，MAE 2.44，polCVaR $-3.81$）に局所解に停滞したのに対し，B1g では3シードすべてで局所解が解消し，M-pol $\CVaR_{0.05}$ も全シードで正値（0.21, 0.89, 1.00，平均 $0.70\pm0.43$）を達成した．

一方，比較対象の B2 リスク感応 IQN は M-est MAE $0.36\pm0.05$（限定 $0.15\pm0.02$）と依然として全手法中最良の精度を示した．
しかし，学習時の $\tau$ サンプリングを歪めた B2p は MAE $3.42\pm1.38$（限定 $5.50\pm3.58$）と，通常の B2 から約10倍悪化した．
B2p の行動選択 polCVaR は $1.01\pm0.02$ と良好な値を維持している（裾の相対順序は保たれている）ものの，絶対的な CVaR 推定精度は大きく損なわれている．

なお M-pol において，B0 は CVaR を過大評価するため3シードすべてで all-RISKY 方策に収束し，収益 CVaR は $-10.79\pm0.00$ となった（$\pm0.00$ となるのは，評価環境の乱数シードが学習シードと独立に固定されているためである）．
B2・B2p は $1.01\pm0.01$ であり，閾値方策クラス内最良値 $1.0000$（$k^*=0$）に最も近い方策性能を示した．
なお，表\ref{tab:B1}で B2 と B2p の M-pol 平均が同値（$1.13\pm0.22$）に見えるのは丸めによる一致であり，per-seed 値は異なる（B2: $1.3809, 1.0000, 1.0000$，B2p: $1.3816, 1.0000, 1.0000$）．

さらに，配置学習の寄与を分離する固定配置の ablation として，FPN を凍結した B3（$\tau_i=i/N$ 固定・算術中点原子）と B3g（$\tau_i=g^{-1}_{\alpha,\lambda}(i/N)$ 固定・$g$-中央値原子，すなわち実験Aの理論配置そのもの）を 3 シードで追加実施した（表\ref{tab:B1}．いずれも FPN 凍結により配置は状態に依存しない）．
B3 の MAE は $4.12\pm0.35$ と B0（7.02）より小さいものの，M-pol は $-7.28\pm1.95$ と all-RISKY 方策に収束しており（一様配置が CVaR を過大評価する挙動は B0 と共通），推定対象の交絡を含む．
B3g は MAE $0.74\pm0.33$・下側裾 15 個（全境界）であり，FPN を学習する B1g（$0.96\pm0.51$）と同程度の推定精度を示した（点推定値は B1g より小さいが，差は 3 シードの標本 SD 内であり優劣は確定できない）．
これは，実験Bにおける提案法の改善の主因が FPN の状態依存最適化そのものではなく，\textbf{歪み分布の理論配置（$g$-一様配置と $g$-中央値原子）の採用}にあることを示唆する．
なお B3g の配置は命題\ref{prop:gmedian}および実験Aの M1 に相当し，提案理論の直接の産物である．B0（$7.02$）・B3（$4.12$）から B3g（$0.74$）への改善は，FPN の有無ではなく\textbf{歪み重み付き量子化の理論そのもの}に帰属する．
toy 環境では収益分布の状態依存性が弱く（all-SAFE 方策下ではほぼ点質量），状態によらずほぼ同一の理論配置が最適であるため最適化の余地が小さいと解釈される．
実際，実験Aでは分布に構造がある設定で M2 が M1 を $6$〜$20$ 倍上回っており（$e_{\mathrm{M2}}/e_{\mathrm{M1}}=0.05$〜$0.17$），分布に状態依存の構造があれば最適化の効果が現れると予想され，状態ごとに最適配置が異なる環境（ヘッジ損益の裾形状が状態に依存する実験C，および toy 環境でもクラッシュ規模を位置依存（例：$-5(1+|T_3|)(1+x/3)$）に変更した中間実験）で検証する（第\ref{sec:future}節）．
なお，方策性能は 3 シードでは優劣が確定できない．CVaR 指標（M-pol CVaR）では B1g が $0.70\pm0.43$ と 3 シードすべて正値で安定する一方，B3g は $0.30\pm0.83$（seed 0 で $-0.66$）とばらつくが，平均収益（M-pol 平均）では B3g の $1.60\pm0.40$ が B1g の $1.37\pm0.36$ を上回る．

\begin{table}[htbp]
\centering
\caption{実験B最終指標（100k ステップ，3シード平均$\pm$SD，SD は標本 SD（ddof=1））．M-est は $\CVaR_{0.05}$ 推定 MAE（参照値は各手法の greedy 方策下の分布に対するものであることに注意．第\ref{subsec:cross}節参照）．「限定」は閾値方策クラス内最良が切り替わる状態に限定した MAE．}
\label{tab:B1}
\small
\setlength{\tabcolsep}{2.5pt}
\begin{tabular}{lccccc}
\toprule
手法 & M-est（全110） & M-est（限定） & M-pol $\CVaR_{0.05}$ & M-pol 平均 & 下側裾 q.f. 数\\
\midrule
B0 FQF & $7.02\pm0.14$ & $14.41\pm0.15$ & $-10.79\pm0.00$ & $3.75\pm0.00$ & $0.00\pm0.00$\\
B1 歪み FQF（一様初期化） & $1.41\pm0.92$ & $3.24\pm2.28$ & $-1.15\pm2.38$ & $2.03\pm0.65$ & $0.48\pm0.46$\\
\textbf{B1g 歪み FQF（$g$-一様初期化）} & $\mathbf{0.96\pm0.51}$ & $\mathbf{0.39\pm0.29}$ & $\mathbf{0.70\pm0.43}$ & $1.37\pm0.36$ & $\mathbf{14.06\pm0.11}$\\
B2 リスク感応 IQN & $0.36\pm0.05$ & $0.15\pm0.02$ & $1.01\pm0.01$ & $1.13\pm0.22$ & （暗黙的表現）\\
B2p IQN（学習時歪み $\tau$） & $3.42\pm1.38$ & $5.50\pm3.58$ & $1.01\pm0.02$ & $1.13\pm0.22$ & （暗黙的表現）\\
\midrule
B3 FQF（一様固定） & $4.12\pm0.35$ & $9.25\pm0.60$ & $-7.28\pm1.95$ & $3.06\pm0.24$ & $0.00\pm0.00$\\
B3g 歪み FQF（$g$-一様固定） & $0.74\pm0.33$ & $0.66\pm0.52$ & $0.30\pm0.83$ & $1.60\pm0.40$ & $15.00\pm0.00$\\
\midrule
閾値方策クラス内最良（$k^*=0$） & --- & --- & $1.0000$ & $1.0000$ & ---\\
\bottomrule
\end{tabular}
\end{table}

\subsubsection{共通方策下の交差評価}
\label{subsec:cross}
M-est の交絡（参照分布が各手法の greedy 方策に依存する）を検討するため，固定方策 $\pi\in\{\text{all-SAFE}（k=0）,\ \text{all-RISKY}（k=10）\}$ の下での参照値 $\widetilde{\mathrm{CVaR}}_{0.05}^{\mathrm{ref}}(s,a;\pi)$ を $10^5$ ロールアウト $\times$ 10 独立バッチのモンテカルロで高精度に計算し，学習済みネットワーク（手法 $\times$ seed の全15個）の CVaR 出力との MAE を全 55 状態 $\times$ 2 行動で比較した（表\ref{tab:B2}）．
バッチ SE（平均 0.011 / 0.035）は，各状態・行動対で得た 10 個の独立バッチ平均値の標本 SD（ddof=1）を $\sqrt{10}$ で割った標準誤差の，全状態・行動対上での平均である．
なお閾値方策 $k=1$ の参照値は，状態・行動レベルの継続方策が $k=0$ と同一である（$k=1$ の差はエピソード先頭の行動選択のみ）ため $k=0$ と一致する．

\begin{table}[htbp]
\centering
\caption{共通方策下の CVaR 推定 MAE（学習済みネットワーク出力 vs 高精度 MC 参照値，55状態$\times$2行動の平均，3シード平均$\pm$SD，ddof=1）．参照値は $10^5$ ロールアウト $\times$ 10 独立バッチ．各ネットワークは学習時の CVaR-greedy 方策を継続方策として学習しているため，参照方策と一致しないセルには継続方策の不一致による差が混入する（本文参照）．}
\label{tab:B2}
\small
\setlength{\tabcolsep}{4pt}
\begin{tabular}{lcc}
\toprule
評価器（学習済み手法） & 参照方策 all-SAFE（$k=0$） & 参照方策 all-RISKY（$k=10$）\\
\midrule
B0 FQF & $1.88\pm0.14$ & $7.01\pm0.14$\\
B1 歪み FQF（一様初期化） & $0.71\pm0.24$ & $5.78\pm0.30$\\
B1g 歪み FQF（$g$-一様初期化） & $0.59\pm0.10$ & $5.52\pm0.18$\\
B2 リスク感応 IQN & $\mathbf{0.35\pm0.05}$ & $5.31\pm0.12$\\
B2p IQN（学習時歪み $\tau$） & $4.19\pm1.48$ & $4.25\pm0.46$\\
\bottomrule
\end{tabular}
\end{table}

表\ref{tab:B2}から次の 2 点が読み取れる．
(i) \textbf{B0 の大きな M-est（7.02）には推定対象の難易度が大きく寄与する}：B0 評価器でも all-SAFE 分布に対する MAE は $1.88$ に低下する一方，all-RISKY（重い裾）参照では全評価器が $4.3$〜$7.0$ と大きな MAE を示す．重い下側裾そのものが主要な難化要因であることが示唆される（ただし評価器間差も残る）．
(ii) \textbf{B2p はいずれの参照に対しても大きな誤差（約4.2）}を示しており，その表現が分布全体として偏っていることを裏付ける（考察の機序説明と整合する）．

ただし，本交差評価には設計上の限界がある．各ネットワークは自身の CVaR-greedy 方策下の収益分布を Bellman 標的として学習しており，critic は固定方策下で再学習していない．したがって表\ref{tab:B2}が測るのは「学習済みネットワーク出力と固定方策参照値の距離（継続方策の不一致を含む）」であり，厳密な固定方策 policy evaluation ではない．参照方策と学習時の継続方策が一致するセル（B0 $\times$ all-RISKY，B1g・B2 $\times$ all-SAFE）以外では，MAE に「表現誤差」だけでなく「継続方策の不一致による真値の差」が混入するため，表\ref{tab:B2}は交絡を完全に分離したものではなく，B0 の $1.88$ を表現誤差と解釈することはできない．
交絡の厳密な分離には，方策を all-RISKY に固定した方策評価モード（行動選択・Bellman 標的の次状態行動・探索をすべて all-RISKY に固定して B0/B1g/B2 を再学習する）での再実験が必要であり，これを第\ref{sec:future}節の課題とする．
現時点で言えるのは，(a) B0 が自身の学習対象（重い裾）に対して $7.01$ の誤差をもつ一方，B1g・B2 は自身の学習対象（ほぼ点質量）に対して $0.59$・$0.35$ にとどまること，(b) 両者の差には推定対象の難易度差が含まれること，の 2 点である．「共通方策下の純粋な表現比較で約3分の1」という読み替えは本設計では正当化されない．

\begin{figure}[htbp]
\centering
\includegraphics[width=0.82\linewidth]{figB2_estimation_ja.png}
\caption{M-est：$\CVaR_{0.05}$ 推定 MAE の学習曲線（5手法，3シード平均$\pm$SD，細線は各シード）．}
\label{fig:B2}
\end{figure}

\begin{figure}[htbp]
\centering
\includegraphics[width=0.88\linewidth]{figB3_fractions_ja.png}
\caption{学習済みquantile fraction配置の分布（B0 一様FQF，B1 一様初期化歪みFQF，B1g $g$-一様初期化歪みFQF）．網掛けは下側裾 $[0,\alpha]$．B1g は全シードで理論配置（14〜15個）を忠実に達成している．}
\label{fig:B3}
\end{figure}

\subsection{考察}
\begin{itemize}
  \item \textbf{原理の検証（実験A）}：真の $\Finv$ の下では提案法はすべての設定で一様配置を大幅に上回った．
  観測された傾き $-1.79$〜$-1.84$ は歪み最適配置 M2 のものであり，一様格子では下側裾の特異性に $O(N^{-(1-1/\nu)})$ で律速される一方，適応配置はそれを回避するという注意\ref{rem:midpoint}の機序と整合する．
  改善比は $N$ 依存で増大するため，予想\ref{conj:alpha}（誤差上界の比 $\Theta(1/\alpha)$）と注意\ref{rem:midpoint}（漸近オーダーの目安 $1/\alpha^2$）を固定 $N$ の実測値と直結させず，傾向の整合として解釈している．
  \item \textbf{RL への移植と初期化（実験B）}：$g$-一様初期化を導入した B1g は，FQF に対して推定誤差を約7分の1（限定状態で約37分の1）に低減した．初期化を理論配置に近づけることで，旧 B1 で見られた seed 0 の局所解トラップが完全に解消され，下側裾への確実な解像度集中が実現された．共通方策下の交差評価（表\ref{tab:B2}）では，all-RISKY 分布（重い裾）の参照で全評価器の MAE が $4.3$〜$7.0$ に達し，推定対象の難易度が誤差を支配することが示唆された．ただし本評価は継続方策の不一致を含むため表現精度だけの分離にはならず，厳密な分離は方策評価モードでの再学習による（第\ref{sec:future}節）．また固定配置 ablation（B3g，MAE $0.74$）が B1g（$0.96$）と同程度の精度を示したことから（差は 3 シードの標本 SD 内），toy 環境での改善の主因は理論配置の採用であり，FPN の状態依存最適化の寄与は本環境では限定的である（表\ref{tab:B1}）．
  \item \textbf{B2p との対照}：B2p の劣化はターゲット側のサンプリングに起因すると解釈される．現在側 level $\tau_i$ の分布を歪めても各 level での分位回帰の一致性は保たれるが，ターゲット側 $\tau'_j\sim g^{-1}(U)$ の等重み標本は $r+\gamma Z'$ ではなく歪み分布 $\dist(r+\gamma Z')$ からの標本となるため，ネットワークは $F^{-1}_{\dist Z}(\tau)=\Finv(g^{-1}(\tau))$ を学習する．読み出し $\tau\in(0,\alpha)$ では $g^{-1}_{\alpha,\lambda}(\tau)\approx\alpha\tau/(1-\lambda)$ という遥かに深い分位点を返し，CVaR を系統的に過小評価する．参照方策に依らず MAE が約 $4.2$ で一定であること（表\ref{tab:B2}）はこの解釈と整合する．重要度重み $1/g'(\tau'_j)$ の付加，現在側のみを歪める変種，読み出し $\tau$ を $g(\tau)$ に写す逆変換変種といった公平な対照は将来課題とする．したがってこの結果は sampling 系の代替法一般の劣位を示すものではなく，また現行の B2p の設定をもって表現側再配分の本質性を実証するものでもない．本実験の範囲での事実は，単純なサンプリング置換では安定した絶対 $\CVaR$ 推定に至らなかった一方，提案法は非裾領域の学習を保持したまま境界と重みを明示的に制御できる点である．
  \item \textbf{IQN との比較と位置づけ}：依然として B2 IQN（0.36）が最良の精度を示しており，提案法の一般的優位性は未確認である．IQN は連続的な分位レベルを入力できる表現力の高さをもつ一方，FQF 系は有限個の離散境界と原子を明示的に保持するため状態ごとの解釈性・可視化性に優れる．本研究の意義は，この離散表現型 RL においてリスク尺度に応じた目的指向量子化を達成した点にある．なお，歪み関数 $g_\alpha$ はリスク水準 $\alpha$ に依存するため，$\alpha$ を変更する場合は FPN の再初期化・再学習を要する．これは学習済み分布を再利用して任意の $\alpha$ で行動選択できる IQN に対する本手法の限界であり，実験C（第\ref{sec:future}節）で定量比較する予定である．
  \item \textbf{有限サンプルの影響}：A-fs で改善幅が縮んだことは，裾にquantile fractionを集中させることが「裾のデータの少なさ」に対する感度と表裏一体であることを示す．
\end{itemize}
%==================== 6 今後の展望 ====================
\section{今後の展望}
\label{sec:future}
\begin{enumerate}
  \item \textbf{デルタヘッジ環境での本実験（実験C）}：Black--Scholes モデル（$K=S_0=100$，$T=30$ 営業日，$\Delta t=1/252$，$\sigma=0.2$，$r=0$ であり，$K$ は行使価格，$S_0$ は初期原資産価格，$T$ は満期，$\Delta t$ は年間252営業日を基準にした日次の時間刻み，$\sigma$ はボラティリティ，$r$ は金利）の下で
  ATM（行使価格＝初期原資産価格）ヨーロピアンコール1枚をショートし，離散ヘッジ比率 $a\in\{0,0.1,\dots,1.0\}$ で日次リヘッジする環境を構築する．
  取引コストは売買想定元本の $0.1\%$，ステップ報酬は日次損益増分，状態は $(S_t/K,$ 満期までの時間$,$ 現ヘッジ比率$)$，
  目的は終端損益の $\CVaR_{0.05}$ 最大化である．
  ベースラインは解析 BS デルタヘッジ（同一コスト・同一離散化），FQF，リスク感応 IQN に加え，深層ヘッジングの先行研究（Cao et al. \cite{cao2021deep,cao2023gamma}）や EX-DRL \cite{malekzadeh2024exdrl} とし，$10^4$ 本のホールドアウトパスで終端損益の経験 $\CVaR_{0.05}$ と平均，
  固定状態での推定 MAE を測る．CVaR 最適化の理論的基盤は Rockafellar and Uryasev \cite{rockafellar2000cvar} に従う．
  さらに，Black--Scholes モデルは連続拡散であり理論的には完全複製可能であるため，裾の重い損益を真にもたらす第一級の拡張として Merton のジャンプ拡散過程や Heston 型の確率的ボラティリティモデルを取り入れた環境での検証を推進する．
  \item \textbf{誤差上界および中点則の定理化}：予想\ref{conj:alpha}の数学的証明を与え，定理として確立することを目指す．
  証明の方針として以下の2ルートを想定している：(i) 中間レポートで示した裾特異性の機序（一様配置は $O(N^{-(1-1/\nu)})$ に律速，適応配置は第1区間幅の急減によりこれを回避）と正則部の中点則二次精度（注意\ref{rem:midpoint}）を厳密な命題として定式化し，符号付き誤差の減衰レートを導出するルート，(ii) Graf and Luschgy \cite{graf2000quantization} の非漸近量子化評価を用いて一様配置の下界と歪み最適配置の上界の比を評価するルートである．
  \item \textbf{ベースライン比較の拡充}：固定配置の ablation（B3/B3g）は実施済みである（第\ref{subsec:expB}節）．残る課題として，(i) 共通方策下の交絡の厳密な分離として，行動選択・Bellman 標的の次状態行動・探索をすべて all-RISKY に固定した\textbf{方策評価モード}で B0/B1g/B2 を再学習し，同一参照と比較する，(ii) B2p の公平な対照として，重要度重み $1/g'(\tau'_j)$ 付加・現在側のみ歪め・読み出しを $g(\tau)$ に写す逆変換の各変種を実装する，(iii) 境界移動と代表点の効果の分離，および Moghimi and Ku \cite{moghimi2025beyond} の静的スペクトルリスク尺度最適化との比較を行う，を予定している．
  toy 環境でもクラッシュ規模を位置依存（例：$-5(1+|T_3|)(1+x/3)$）に変更すれば状態依存性を低コストで導入でき，FPN 学習の寄与を数時間規模で先行的に検証できる．
  状態ごとに最適配置が異なる環境（実験C のヘッジ損益）では FPN の状態依存最適化が効くかを検証する．
  実験Bはシード数を 10 以上に拡張して統計的な優位性（bootstrap 信頼区間，paired comparison を含む）を確認するとともに，
  独立な複数評価バッチから参照値の標準誤差を評価する．
  \item \textbf{時間整合性}：静的 CVaR 目的に対して時間整合な方策クラス（Lim and Malik \cite{lim2022distributional}）との組合せを検討し，
  推定精度の改善を方策性能の改善に接続する．
\end{enumerate}
%==================== 参考文献 ====================
\addcontentsline{toc}{section}{参考文献}
\begin{thebibliography}{99}
\small
\bibitem{yang2019fqf} D. Yang, L. Zhao, Z. Lin, T. Qin, J. Bian, T.-Y. Liu.
Fully Parameterized Quantile Function for Distributional Reinforcement Learning. \textit{NeurIPS 2019}.
\bibitem{dabney2018iqn} W. Dabney, G. Ostrovski, D. Silver, R. Munos.
Implicit Quantile Networks for Distributional Reinforcement Learning. \textit{ICML 2018}.
\bibitem{dabney2018qrdqn} W. Dabney, M. Rowland, M. G. Bellemare, R. Munos.
Distributional Reinforcement Learning with Quantile Regression. \textit{AAAI 2018}.
\bibitem{bellemare2017c51} M. G. Bellemare, W. Dabney, R. Munos.
A Distributional Perspective on Reinforcement Learning. \textit{ICML 2017}.
\bibitem{moghimi2025beyond} M. Moghimi, H. Ku.
Beyond CVaR: Leveraging Static Spectral Risk Measures for Enhanced Decision-Making in Distributional Reinforcement Learning. \textit{ICML 2025}.
\bibitem{zhang2025tailsafe} J. Zhang.
Tail-Safe Hedging: Explainable Risk-Sensitive Reinforcement Learning with a White-Box CBF--QP Safety Layer in Arbitrage-Free Markets.
arXiv:2510.04555, 2025.（単著・査読なしプレプリント）
\bibitem{graf2000quantization} S. Graf, H. Luschgy.
Foundations of Quantization for Probability Distributions. \textit{Lecture Notes in Mathematics} 1730, Springer, 2000.
\bibitem{rockafellar2000cvar} R. T. Rockafellar, S. Uryasev.
Optimization of Conditional Value-at-Risk. \textit{Journal of Risk} 2(3): 21--41, 2000.
\bibitem{lim2022distributional} S. H. Lim, I. Malik.
Distributional Reinforcement Learning for Risk-Sensitive Policies. \textit{NeurIPS 2022}.
\bibitem{yaari1987dual} M. E. Yaari.
The Dual Theory of Choice under Risk. \textit{Econometrica} 55(1): 95--115, 1987.
\bibitem{pichler2013evaluations} A. Pichler.
Evaluations of Risk Measures for Different Probability Measures. \textit{SIAM Journal on Optimization} 23(1): 530--551, 2013.
\bibitem{kiesel2016wasserstein} R. Kiesel, R. R{\"u}hlicke, G. Stahl, J. Zheng.
The Wasserstein Metric and Robustness in Risk Management. \textit{Risks} 4(3): 32, 2016.
\bibitem{gan2026energy} G. Gan, N. Chen, Z. Wei, H. Qin, Y. Lu, H. Huang, Y. Huang.
Energy-aware job scheduling for green data centers under uncertainty: A structure-compressed spectral risk control approach. \textit{Applied Energy} 420: 128137, 2026.
\bibitem{pflug2014multistage} G. C. Pflug, A. Pichler.
\textit{Multistage Stochastic Optimization}. Springer Series in Operations Research and Financial Engineering, Springer, 2014.
\bibitem{cao2021deep} J. Cao, J. Chen, J. Hull, Z. Poulos.
Deep Hedging of Derivatives Using Reinforcement Learning. \textit{The Journal of Financial Data Science} 3(1): 10--27, 2021.
\bibitem{cao2023gamma} J. Cao, J. Chen, S. Farghadani, J. Hull, Z. Poulos, Z. Wang, J. Yuan.
Gamma and vega hedging using deep distributional reinforcement learning. \textit{Frontiers in Artificial Intelligence} 6: 1129370, 2023.
\bibitem{malekzadeh2024exdrl} P. Malekzadeh, Z. Poulos, J. Chen, Z. Wang, K. N. Plataniotis.
EX-DRL: Hedging Against Heavy Losses with EXtreme Distributional Reinforcement Learning. arXiv:2408.12446, 2024.（プレプリント）
\bibitem{zhang2026rqiqn} Z. Zhang, M. Yang, R. Chen, S. Xie, H. Xiong.
Quantile Geometry Regularization for Distributional Reinforcement Learning. arXiv:2605.08182, 2026.（プレプリント）
\end{thebibliography}
\end{document}
